% !TEX program = lualatex
% ---------------------------------------------------------------------------------------------
%  Readback note 45: Cup products of crystalline classes vanish up to a power of p: the walk along Frobenius orbits
%
%  Build:  python3 tools/rb.py note <slug>
%
%  One unit of a Lean artifact, rewritten as ordinary mathematics. The prose is hand-written. Lean names,
%  line ranges and counts come from data.tex; quoted Lean statements are copied from the artifact by the build.
%
%  Rules of the text: standard terminology; every symbol is explained where it first appears and listed again
%  in the notation at the end; no letter has two meanings; a proof is complete or says what it takes for granted;
%  Lean names appear only in the LEAN lines and in the concordance.
%
%  The signature rule (the build enforces it): a function, relation or predicate that the artifact defines is
%  written only through its notation macro, which the build makes from the spec (`notation`) with one slot for
%  each explicit argument of the Lean definition. No argument is fixed and dropped, also not "for brevity" and
%  not for one section. Every argument and every hypothesis of a stated declaration stands inside its
%  statement, or in a Setting that the statement cites by number; only implicit and instance arguments (the
%  ambient objects) may stand in the conventions.
% ---------------------------------------------------------------------------------------------
\documentclass[11pt,a4paper]{article}
\ifdefined\pdfvariable\pdfvariable suppressoptionalinfo \numexpr 1+2+4+8\relax\fi
\usepackage[a4paper,hmargin=35mm,top=22mm,bottom=27mm,footskip=12mm,nohead]{geometry}

\usepackage{fontspec}
\directlua{
  luaotfload.add_fallback("gfallback", {"STIXTwoMath-Regular.otf:mode=node;", "DejaVuSans.ttf:mode=node;"})
  luaotfload.add_fallback("gmonofallback", {"DejaVuSansMono.ttf:mode=node;", "STIXTwoMath-Regular.otf:mode=node;"})
}
\setmainfont{LibertinusSerif}[Extension=.otf,UprightFont=*-Regular,ItalicFont=*-Italic,
  BoldFont=*-Semibold,BoldItalicFont=*-SemiboldItalic,RawFeature={fallback=gfallback}]
\setsansfont{Inter}[Extension=.otf,UprightFont=*-Regular,ItalicFont=*-Italic,
  BoldFont=*-SemiBold,BoldItalicFont=*-SemiBoldItalic,RawFeature={fallback=gfallback}]
\setmonofont{JetBrainsMono}[Extension=.otf,UprightFont=*-Regular,ItalicFont=*-Italic,
  BoldFont=*-SemiBold,BoldItalicFont=*-SemiBoldItalic,Scale=MatchLowercase,
  RawFeature={-calt;-liga;fallback=gmonofallback}]
\usepackage{amsmath}
\usepackage{amsthm}
\usepackage[math-style=TeX]{unicode-math}
\setmathfont{LibertinusMath-Regular.otf}

\usepackage{xcolor}
\usepackage{tikz}
\usetikzlibrary{calc,arrows.meta,positioning,backgrounds}
\usepackage[protrusion=true,expansion=true]{microtype}
\usepackage{enumitem}
\usepackage{needspace}
\usepackage{booktabs}
\usepackage{array}
\usepackage{longtable}
\usepackage{fancyhdr}
\usepackage{lastpage}
\usepackage{multicol}
\usepackage[hidelinks,pdfusetitle]{hyperref}

\definecolor{ink}{HTML}{0F172A}
\definecolor{inksoft}{HTML}{475569}
\definecolor{inkfaint}{HTML}{94A3B8}
\definecolor{rulegrey}{HTML}{CBD5E1}
\definecolor{arith}{HTML}{C2410C}
\definecolor{est}{HTML}{7C3AED}
\definecolor{ana}{HTML}{0891B2}
\definecolor{ok}{HTML}{047857}
\definecolor{leancol}{HTML}{3730A3}

\title{Cup products of crystalline classes vanish up to a power of p: the walk along Frobenius orbits}
\author{grh-compact}
\date{}

% ------------------------------------------------------------------------------------- macros
\newcommand{\leanfont}{\ttfamily\fontsize{8.6}{10.5}\selectfont\color{leancol}}
\newcommand{\lean}[1]{\csname LN@#1\endcsname}
\newcommand{\fillin}[1]{{\color{arith}\itshape[#1]}}  % a place of the template that is not written yet; the build reports every one that is left
\newcommand{\leanquote}[1]{\csname LQ@#1\endcsname}
% generated by rb.py note; do not edit
\newcommand{\ConcDecls}{44}
\newcommand{\Pin}{72fead10}
\newcommand{\PinText}{that of commit {\leanfont 72fead10}}
\newcommand{\ArtifactName}{grh-compact}
\expandafter\newcommand\csname LN@wittFst\endcsname{\href{https://github.com/lsernau/grh-lean/blob/72fead102f663d5c7a407964a13dc0cf9d977207/grh-compact/src/New/Lift/CupProductBound.lean\#L40-L40}{{\leanfont Lift.\allowbreak{}KummerChart.\allowbreak{}PairingSide.\allowbreak{}wittFst}}}
\expandafter\newcommand\csname LN@tensorLattice\endcsname{\href{https://github.com/lsernau/grh-lean/blob/72fead102f663d5c7a407964a13dc0cf9d977207/grh-compact/src/New/Lift/CupProductBound.lean\#L44-L45}{{\leanfont Lift.\allowbreak{}KummerChart.\allowbreak{}PairingSide.\allowbreak{}tensorLattice}}}
\expandafter\newcommand\csname LN@thirdSide\endcsname{\href{https://github.com/lsernau/grh-lean/blob/72fead102f663d5c7a407964a13dc0cf9d977207/grh-compact/src/New/Lift/CupProductBound.lean\#L51-L53}{{\leanfont Lift.\allowbreak{}Pairing.\allowbreak{}thirdSide}}}
\expandafter\newcommand\csname LN@wittPairing\endcsname{\href{https://github.com/lsernau/grh-lean/blob/72fead102f663d5c7a407964a13dc0cf9d977207/grh-compact/src/New/Lift/CupProductBound.lean\#L57-L58}{{\leanfont Lift.\allowbreak{}Pairing.\allowbreak{}wittPairing}}}
\expandafter\newcommand\csname LN@IsSafe\endcsname{\href{https://github.com/lsernau/grh-lean/blob/72fead102f663d5c7a407964a13dc0cf9d977207/grh-compact/src/New/Lift/CupProductBound.lean\#L62-L62}{{\leanfont Lift.\allowbreak{}Pairing.\allowbreak{}IsSafe}}}
\expandafter\newcommand\csname LN@ScaledCupBounds\endcsname{\href{https://github.com/lsernau/grh-lean/blob/72fead102f663d5c7a407964a13dc0cf9d977207/grh-compact/src/New/Lift/CupProductBound.lean\#L66-L67}{{\leanfont Lift.\allowbreak{}Pairing.\allowbreak{}ScaledCupBounds}}}
\expandafter\newcommand\csname LN@wittFstsmul\endcsname{\href{https://github.com/lsernau/grh-lean/blob/72fead102f663d5c7a407964a13dc0cf9d977207/grh-compact/src/New/Lift/CupProductBound.lean\#L69-L71}{{\leanfont Lift.\allowbreak{}wittFst\_smul}}}
\expandafter\newcommand\csname LN@smulwittPairing\endcsname{\href{https://github.com/lsernau/grh-lean/blob/72fead102f663d5c7a407964a13dc0cf9d977207/grh-compact/src/New/Lift/CupProductBound.lean\#L73-L74}{{\leanfont Lift.\allowbreak{}Pairing.\allowbreak{}smul\_wittPairing}}}
\expandafter\newcommand\csname LN@ofisSafe\endcsname{\href{https://github.com/lsernau/grh-lean/blob/72fead102f663d5c7a407964a13dc0cf9d977207/grh-compact/src/New/Lift/CupProductBound.lean\#L83-L122}{{\leanfont Lift.\allowbreak{}Pairing.\allowbreak{}ScaledCupBounds.\allowbreak{}of\_isSafe}}}
\expandafter\newcommand\csname LQ@ofisSafe\endcsname{\leanquoteline{0}{lemma Pairing.ScaledCupBounds.of\_isSafe \{e₁ e₂ e₃ c c'\} x w x' (he : e₁ + e₂ + e₃ ≤ e) (hx :}\leanquoteline{2}{S.projTriv.wittMap x = p \textasciicircum{} e₁ • .mk 1 c) (hw : S.projTriv.wittMap w = p \textasciicircum{} e₁ • .mk 1 c) (hx' :}\leanquoteline{2}{S'.projTriv.wittMap x' = p \textasciicircum{} e₂ • .mk 1 c') (H : ∀ g : k.galMonoid, (∀ h :}\leanquoteline{2}{k.galMonoid, P.IsSafe e₃ (g • x - x) (g • (h • x' - x'))) ∧ P.IsSafe e₃ (w - x) (g • x' - x') ∧}\leanquoteline{2}{P.IsSafe e₃ (g • w - w) (g • (φ • x' - x')) ∧ P.IsSafe e₃ (φ • x - w) (g • φ • x' - φ • x')) :}\leanquoteline{2}{P.ScaledCupBounds e c c'}}
\expandafter\newcommand\csname LN@tateLattice\endcsname{\href{https://github.com/lsernau/grh-lean/blob/72fead102f663d5c7a407964a13dc0cf9d977207/grh-compact/src/New/Lift/LongGap.lean\#L25-L27}{{\leanfont Lift.\allowbreak{}KummerChart.\allowbreak{}tateLattice}}}
\expandafter\newcommand\csname LN@stepBackzero\endcsname{\href{https://github.com/lsernau/grh-lean/blob/72fead102f663d5c7a407964a13dc0cf9d977207/grh-compact/src/New/Lift/LongGap.lean\#L37-L37}{{\leanfont Lift.\allowbreak{}KummerChart.\allowbreak{}PairingSide.\allowbreak{}stepBack\_zero}}}
\expandafter\newcommand\csname LN@inclKer\endcsname{\href{https://github.com/lsernau/grh-lean/blob/72fead102f663d5c7a407964a13dc0cf9d977207/grh-compact/src/New/Lift/LongGap.lean\#L40-L40}{{\leanfont Lift.\allowbreak{}KummerChart.\allowbreak{}PairingSide.\allowbreak{}inclKer}}}
\expandafter\newcommand\csname LN@projTrivLoc\endcsname{\href{https://github.com/lsernau/grh-lean/blob/72fead102f663d5c7a407964a13dc0cf9d977207/grh-compact/src/New/Lift/LongGap.lean\#L44-L45}{{\leanfont Lift.\allowbreak{}KummerChart.\allowbreak{}PairingSide.\allowbreak{}projTrivLoc}}}
\expandafter\newcommand\csname LN@deepElements\endcsname{\href{https://github.com/lsernau/grh-lean/blob/72fead102f663d5c7a407964a13dc0cf9d977207/grh-compact/src/New/Lift/LongGap.lean\#L50-L51}{{\leanfont Lift.\allowbreak{}KummerChart.\allowbreak{}PairingSide.\allowbreak{}deepElements}}}
\expandafter\newcommand\csname LN@kerWithExponents\endcsname{\href{https://github.com/lsernau/grh-lean/blob/72fead102f663d5c7a407964a13dc0cf9d977207/grh-compact/src/New/Lift/LongGap.lean\#L55-L56}{{\leanfont Lift.\allowbreak{}KummerChart.\allowbreak{}PairingSide.\allowbreak{}kerWithExponents}}}
\expandafter\newcommand\csname LN@Reaches\endcsname{\href{https://github.com/lsernau/grh-lean/blob/72fead102f663d5c7a407964a13dc0cf9d977207/grh-compact/src/New/Lift/LongGap.lean\#L61-L62}{{\leanfont Lift.\allowbreak{}KummerChart.\allowbreak{}PairingSide.\allowbreak{}Reaches}}}
\expandafter\newcommand\csname LN@SidePairing\endcsname{\href{https://github.com/lsernau/grh-lean/blob/72fead102f663d5c7a407964a13dc0cf9d977207/grh-compact/src/New/Lift/LongGap.lean\#L65-L65}{{\leanfont Lift.\allowbreak{}KummerChart.\allowbreak{}PairingSide.\allowbreak{}SidePairing}}}
\expandafter\newcommand\csname LN@pointsOver\endcsname{\href{https://github.com/lsernau/grh-lean/blob/72fead102f663d5c7a407964a13dc0cf9d977207/grh-compact/src/New/Lift/LongGap.lean\#L69-L70}{{\leanfont Lift.\allowbreak{}KummerChart.\allowbreak{}PairingSide.\allowbreak{}pointsOver}}}
\expandafter\newcommand\csname LN@HasDeepPoint\endcsname{\href{https://github.com/lsernau/grh-lean/blob/72fead102f663d5c7a407964a13dc0cf9d977207/grh-compact/src/New/Lift/LongGap.lean\#L75-L79}{{\leanfont Lift.\allowbreak{}KummerChart.\allowbreak{}PairingSide.\allowbreak{}HasDeepPoint}}}
\expandafter\newcommand\csname LQ@HasDeepPoint\endcsname{\leanquoteline{0}{def HasDeepPoint (cut : \_ → Prop) l (o o' z : \_ → ℕ → ℤ → ℤ) (Q : Prop) := ∃ d r,}\leanquoteline{4}{(S.pointsOver c' i' e (o i') d).Nonempty ∧ cut (i'.stepBack r) ∧ ∀ a, ∀ u ∈ S.U (i'.stepBack a),}\leanquoteline{2}{∀ v ∈ S'.U (i'.stepBack a), -z i' a (-u - v - 1) - d a u < min (-v) (o' i' a v) ∧}\leanquoteline{2}{(Q → stdExponent l a u ≤ o i' a u ∧ -frobExponent r 0 (z (i'.stepBack r)) a (-u - v - 1) - d a u <}\leanquoteline{2}{min (o' i' a v) (frobExponent l 0 (o' (i'.stepBack l)) a v))}}
\expandafter\newcommand\csname LN@wittFstinclKer\endcsname{\href{https://github.com/lsernau/grh-lean/blob/72fead102f663d5c7a407964a13dc0cf9d977207/grh-compact/src/New/Lift/LongGap.lean\#L88-L89}{{\leanfont Lift.\allowbreak{}KummerChart.\allowbreak{}PairingSide.\allowbreak{}wittFst\_inclKer}}}
\expandafter\newcommand\csname LN@isSafesum\endcsname{\href{https://github.com/lsernau/grh-lean/blob/72fead102f663d5c7a407964a13dc0cf9d977207/grh-compact/src/New/Lift/LongGap.lean\#L91-L100}{{\leanfont Lift.\allowbreak{}KummerChart.\allowbreak{}PairingSide.\allowbreak{}isSafe\_sum}}}
\expandafter\newcommand\csname LN@existsextendform\endcsname{\href{https://github.com/lsernau/grh-lean/blob/72fead102f663d5c7a407964a13dc0cf9d977207/grh-compact/src/New/Lift/LongGap.lean\#L105-L120}{{\leanfont Lift.\allowbreak{}KummerChart.\allowbreak{}PairingSide.\allowbreak{}exists\_extend\_form}}}
\expandafter\newcommand\csname LN@memofstable\endcsname{\href{https://github.com/lsernau/grh-lean/blob/72fead102f663d5c7a407964a13dc0cf9d977207/grh-compact/src/New/Lift/LongGap.lean\#L122-L124}{{\leanfont Lift.\allowbreak{}KummerChart.\allowbreak{}PairingSide.\allowbreak{}Reaches.\allowbreak{}mem\_of\_stable}}}
\expandafter\newcommand\csname LN@kerWithExponentsledeepElements\endcsname{\href{https://github.com/lsernau/grh-lean/blob/72fead102f663d5c7a407964a13dc0cf9d977207/grh-compact/src/New/Lift/LongGap.lean\#L126-L127}{{\leanfont Lift.\allowbreak{}KummerChart.\allowbreak{}PairingSide.\allowbreak{}kerWithExponents\_le\_deepElements}}}
\expandafter\newcommand\csname LN@existsinclKereqofprojTriveqzero\endcsname{\href{https://github.com/lsernau/grh-lean/blob/72fead102f663d5c7a407964a13dc0cf9d977207/grh-compact/src/New/Lift/LongGap.lean\#L132-L148}{{\leanfont Lift.\allowbreak{}KummerChart.\allowbreak{}PairingSide.\allowbreak{}exists\_inclKer\_eq\_of\_projTriv\_eq\_zero}}}
\expandafter\newcommand\csname LN@pointsOvernonemptyofkillsCokernel\endcsname{\href{https://github.com/lsernau/grh-lean/blob/72fead102f663d5c7a407964a13dc0cf9d977207/grh-compact/src/New/Lift/LongGap.lean\#L153-L159}{{\leanfont Lift.\allowbreak{}KummerChart.\allowbreak{}PairingSide.\allowbreak{}pointsOver\_nonempty\_of\_killsCokernel}}}
\expandafter\newcommand\csname LN@existsreachesdeep\endcsname{\href{https://github.com/lsernau/grh-lean/blob/72fead102f663d5c7a407964a13dc0cf9d977207/grh-compact/src/New/Lift/LongGap.lean\#L165-L191}{{\leanfont Lift.\allowbreak{}KummerChart.\allowbreak{}PairingSide.\allowbreak{}exists\_reaches\_deep}}}
\expandafter\newcommand\csname LN@pointsOvernonemptyofcosts\endcsname{\href{https://github.com/lsernau/grh-lean/blob/72fead102f663d5c7a407964a13dc0cf9d977207/grh-compact/src/New/Lift/LongGap.lean\#L197-L206}{{\leanfont Lift.\allowbreak{}KummerChart.\allowbreak{}PairingSide.\allowbreak{}pointsOver\_nonempty\_of\_costs}}}
\expandafter\newcommand\csname LN@smulsubmemkerWithExponents\endcsname{\href{https://github.com/lsernau/grh-lean/blob/72fead102f663d5c7a407964a13dc0cf9d977207/grh-compact/src/New/Lift/LongGap.lean\#L208-L211}{{\leanfont Lift.\allowbreak{}KummerChart.\allowbreak{}PairingSide.\allowbreak{}smul\_sub\_mem\_kerWithExponents}}}
\expandafter\newcommand\csname LN@submemkerWithExponentsofpointsOver\endcsname{\href{https://github.com/lsernau/grh-lean/blob/72fead102f663d5c7a407964a13dc0cf9d977207/grh-compact/src/New/Lift/LongGap.lean\#L213-L216}{{\leanfont Lift.\allowbreak{}KummerChart.\allowbreak{}PairingSide.\allowbreak{}sub\_mem\_kerWithExponents\_of\_pointsOver}}}
\expandafter\newcommand\csname LN@phipowsmulmempointsOver\endcsname{\href{https://github.com/lsernau/grh-lean/blob/72fead102f663d5c7a407964a13dc0cf9d977207/grh-compact/src/New/Lift/LongGap.lean\#L218-L227}{{\leanfont Lift.\allowbreak{}KummerChart.\allowbreak{}PairingSide.\allowbreak{}phi\_pow\_smul\_mem\_pointsOver}}}
\expandafter\newcommand\csname LN@wittPairingconstantAtandthirdPairing\endcsname{\href{https://github.com/lsernau/grh-lean/blob/72fead102f663d5c7a407964a13dc0cf9d977207/grh-compact/src/New/Lift/LongGap.lean\#L229-L237}{{\leanfont Lift.\allowbreak{}KummerChart.\allowbreak{}PairingSide.\allowbreak{}wittPairing\_constantAt\_and\_thirdPairing}}}
\expandafter\newcommand\csname LN@isSafeofexponents\endcsname{\href{https://github.com/lsernau/grh-lean/blob/72fead102f663d5c7a407964a13dc0cf9d977207/grh-compact/src/New/Lift/LongGap.lean\#L244-L259}{{\leanfont Lift.\allowbreak{}KummerChart.\allowbreak{}PairingSide.\allowbreak{}SidePairing.\allowbreak{}isSafe\_of\_exponents}}}
\expandafter\newcommand\csname LN@isSafesmulsub\endcsname{\href{https://github.com/lsernau/grh-lean/blob/72fead102f663d5c7a407964a13dc0cf9d977207/grh-compact/src/New/Lift/LongGap.lean\#L265-L267}{{\leanfont Lift.\allowbreak{}KummerChart.\allowbreak{}PairingSide.\allowbreak{}SidePairing.\allowbreak{}isSafe\_smul\_sub}}}
\expandafter\newcommand\csname LN@scaledCupBoundsofpointsatcuts\endcsname{\href{https://github.com/lsernau/grh-lean/blob/72fead102f663d5c7a407964a13dc0cf9d977207/grh-compact/src/New/Lift/LongGap.lean\#L273-L309}{{\leanfont Lift.\allowbreak{}KummerChart.\allowbreak{}PairingSide.\allowbreak{}SidePairing.\allowbreak{}scaledCupBounds\_of\_points\_at\_cuts}}}
\expandafter\newcommand\csname LN@pointsOvernonemptyofnoMultipleNear\endcsname{\href{https://github.com/lsernau/grh-lean/blob/72fead102f663d5c7a407964a13dc0cf9d977207/grh-compact/src/New/Lift/LongGap.lean\#L317-L323}{{\leanfont Lift.\allowbreak{}KummerChart.\allowbreak{}PairingSide.\allowbreak{}pointsOver\_nonempty\_of\_noMultipleNear}}}
\expandafter\newcommand\csname LN@scaledCupBoundsofkillsCokernel\endcsname{\href{https://github.com/lsernau/grh-lean/blob/72fead102f663d5c7a407964a13dc0cf9d977207/grh-compact/src/New/Lift/LongGap.lean\#L331-L350}{{\leanfont Lift.\allowbreak{}KummerChart.\allowbreak{}PairingSide.\allowbreak{}SidePairing.\allowbreak{}scaledCupBounds\_of\_killsCokernel}}}
\expandafter\newcommand\csname LN@raiseExponent\endcsname{\href{https://github.com/lsernau/grh-lean/blob/72fead102f663d5c7a407964a13dc0cf9d977207/grh-compact/src/New/Lift/LongGap.lean\#L359-L361}{{\leanfont Lift.\allowbreak{}KummerChart.\allowbreak{}JumpLattice.\allowbreak{}raiseExponent}}}
\expandafter\newcommand\csname LN@existscombofjumpMultiplicityle\endcsname{\href{https://github.com/lsernau/grh-lean/blob/72fead102f663d5c7a407964a13dc0cf9d977207/grh-compact/src/New/Lift/LongGap.lean\#L366-L404}{{\leanfont Lift.\allowbreak{}KummerChart.\allowbreak{}JumpLattice.\allowbreak{}exists\_comb\_of\_jumpMultiplicity\_le}}}
\expandafter\newcommand\csname LN@shiftedStd\endcsname{\href{https://github.com/lsernau/grh-lean/blob/72fead102f663d5c7a407964a13dc0cf9d977207/grh-compact/src/New/Lift/LongGap.lean\#L417-L417}{{\leanfont Lift.\allowbreak{}KummerChart.\allowbreak{}JumpLattice.\allowbreak{}shiftedStd}}}
\expandafter\newcommand\csname LN@constantsModE\endcsname{\href{https://github.com/lsernau/grh-lean/blob/72fead102f663d5c7a407964a13dc0cf9d977207/grh-compact/src/New/Lift/LongGap.lean\#L421-L422}{{\leanfont Lift.\allowbreak{}KummerChart.\allowbreak{}JumpLattice.\allowbreak{}constantsModE}}}
\expandafter\newcommand\csname LN@coeffmemconstantsModE\endcsname{\href{https://github.com/lsernau/grh-lean/blob/72fead102f663d5c7a407964a13dc0cf9d977207/grh-compact/src/New/Lift/LongGap.lean\#L430-L448}{{\leanfont Lift.\allowbreak{}KummerChart.\allowbreak{}JumpLattice.\allowbreak{}coeff\_mem\_constantsModE}}}
\expandafter\newcommand\csname LN@letterProjmemconstantsModE\endcsname{\href{https://github.com/lsernau/grh-lean/blob/72fead102f663d5c7a407964a13dc0cf9d977207/grh-compact/src/New/Lift/LongGap.lean\#L453-L458}{{\leanfont Lift.\allowbreak{}KummerChart.\allowbreak{}JumpLattice.\allowbreak{}letterProj\_mem\_constantsModE}}}
\expandafter\newcommand\csname LN@pointsOvernonemptyoflonggap\endcsname{\href{https://github.com/lsernau/grh-lean/blob/72fead102f663d5c7a407964a13dc0cf9d977207/grh-compact/src/New/Lift/LongGap.lean\#L463-L496}{{\leanfont Lift.\allowbreak{}KummerChart.\allowbreak{}PairingSide.\allowbreak{}pointsOver\_nonempty\_of\_long\_gap}}}
\newcommand{\ConcordanceRows}{\crow{wittFst}{\href{https://github.com/lsernau/grh-lean/blob/72fead102f663d5c7a407964a13dc0cf9d977207/grh-compact/src/New/Lift/CupProductBound.lean\#L40-L40}{{\leanfont Lift.\allowbreak{}KummerChart.\allowbreak{}PairingSide.\allowbreak{}wittFst}}}{CupProductBound.lean}{40--40}
\crow{tensorLattice}{\href{https://github.com/lsernau/grh-lean/blob/72fead102f663d5c7a407964a13dc0cf9d977207/grh-compact/src/New/Lift/CupProductBound.lean\#L44-L45}{{\leanfont Lift.\allowbreak{}KummerChart.\allowbreak{}PairingSide.\allowbreak{}tensorLattice}}}{CupProductBound.lean}{44--45}
\crow{thirdSide}{\href{https://github.com/lsernau/grh-lean/blob/72fead102f663d5c7a407964a13dc0cf9d977207/grh-compact/src/New/Lift/CupProductBound.lean\#L51-L53}{{\leanfont Lift.\allowbreak{}Pairing.\allowbreak{}thirdSide}}}{CupProductBound.lean}{51--53}
\crow{wittPairing}{\href{https://github.com/lsernau/grh-lean/blob/72fead102f663d5c7a407964a13dc0cf9d977207/grh-compact/src/New/Lift/CupProductBound.lean\#L57-L58}{{\leanfont Lift.\allowbreak{}Pairing.\allowbreak{}wittPairing}}}{CupProductBound.lean}{57--58}
\crow{IsSafe}{\href{https://github.com/lsernau/grh-lean/blob/72fead102f663d5c7a407964a13dc0cf9d977207/grh-compact/src/New/Lift/CupProductBound.lean\#L62-L62}{{\leanfont Lift.\allowbreak{}Pairing.\allowbreak{}IsSafe}}}{CupProductBound.lean}{62--62}
\crow{ScaledCupBounds}{\href{https://github.com/lsernau/grh-lean/blob/72fead102f663d5c7a407964a13dc0cf9d977207/grh-compact/src/New/Lift/CupProductBound.lean\#L66-L67}{{\leanfont Lift.\allowbreak{}Pairing.\allowbreak{}ScaledCupBounds}}}{CupProductBound.lean}{66--67}
\crow{wittFstsmul}{\href{https://github.com/lsernau/grh-lean/blob/72fead102f663d5c7a407964a13dc0cf9d977207/grh-compact/src/New/Lift/CupProductBound.lean\#L69-L71}{{\leanfont Lift.\allowbreak{}wittFst\_smul}}}{CupProductBound.lean}{69--71}
\crow{smulwittPairing}{\href{https://github.com/lsernau/grh-lean/blob/72fead102f663d5c7a407964a13dc0cf9d977207/grh-compact/src/New/Lift/CupProductBound.lean\#L73-L74}{{\leanfont Lift.\allowbreak{}Pairing.\allowbreak{}smul\_wittPairing}}}{CupProductBound.lean}{73--74}
\crow{ofisSafe}{\href{https://github.com/lsernau/grh-lean/blob/72fead102f663d5c7a407964a13dc0cf9d977207/grh-compact/src/New/Lift/CupProductBound.lean\#L83-L122}{{\leanfont Lift.\allowbreak{}Pairing.\allowbreak{}ScaledCupBounds.\allowbreak{}of\_isSafe}}}{CupProductBound.lean}{83--122}
\crow{tateLattice}{\href{https://github.com/lsernau/grh-lean/blob/72fead102f663d5c7a407964a13dc0cf9d977207/grh-compact/src/New/Lift/LongGap.lean\#L25-L27}{{\leanfont Lift.\allowbreak{}KummerChart.\allowbreak{}tateLattice}}}{LongGap.lean}{25--27}
\crow{stepBackzero}{\href{https://github.com/lsernau/grh-lean/blob/72fead102f663d5c7a407964a13dc0cf9d977207/grh-compact/src/New/Lift/LongGap.lean\#L37-L37}{{\leanfont Lift.\allowbreak{}KummerChart.\allowbreak{}PairingSide.\allowbreak{}stepBack\_zero}}}{LongGap.lean}{37--37}
\crow{inclKer}{\href{https://github.com/lsernau/grh-lean/blob/72fead102f663d5c7a407964a13dc0cf9d977207/grh-compact/src/New/Lift/LongGap.lean\#L40-L40}{{\leanfont Lift.\allowbreak{}KummerChart.\allowbreak{}PairingSide.\allowbreak{}inclKer}}}{LongGap.lean}{40--40}
\crow{projTrivLoc}{\href{https://github.com/lsernau/grh-lean/blob/72fead102f663d5c7a407964a13dc0cf9d977207/grh-compact/src/New/Lift/LongGap.lean\#L44-L45}{{\leanfont Lift.\allowbreak{}KummerChart.\allowbreak{}PairingSide.\allowbreak{}projTrivLoc}}}{LongGap.lean}{44--45}
\crow{deepElements}{\href{https://github.com/lsernau/grh-lean/blob/72fead102f663d5c7a407964a13dc0cf9d977207/grh-compact/src/New/Lift/LongGap.lean\#L50-L51}{{\leanfont Lift.\allowbreak{}KummerChart.\allowbreak{}PairingSide.\allowbreak{}deepElements}}}{LongGap.lean}{50--51}
\crow{kerWithExponents}{\href{https://github.com/lsernau/grh-lean/blob/72fead102f663d5c7a407964a13dc0cf9d977207/grh-compact/src/New/Lift/LongGap.lean\#L55-L56}{{\leanfont Lift.\allowbreak{}KummerChart.\allowbreak{}PairingSide.\allowbreak{}kerWithExponents}}}{LongGap.lean}{55--56}
\crow{Reaches}{\href{https://github.com/lsernau/grh-lean/blob/72fead102f663d5c7a407964a13dc0cf9d977207/grh-compact/src/New/Lift/LongGap.lean\#L61-L62}{{\leanfont Lift.\allowbreak{}KummerChart.\allowbreak{}PairingSide.\allowbreak{}Reaches}}}{LongGap.lean}{61--62}
\crow{SidePairing}{\href{https://github.com/lsernau/grh-lean/blob/72fead102f663d5c7a407964a13dc0cf9d977207/grh-compact/src/New/Lift/LongGap.lean\#L65-L65}{{\leanfont Lift.\allowbreak{}KummerChart.\allowbreak{}PairingSide.\allowbreak{}SidePairing}}}{LongGap.lean}{65--65}
\crow{pointsOver}{\href{https://github.com/lsernau/grh-lean/blob/72fead102f663d5c7a407964a13dc0cf9d977207/grh-compact/src/New/Lift/LongGap.lean\#L69-L70}{{\leanfont Lift.\allowbreak{}KummerChart.\allowbreak{}PairingSide.\allowbreak{}pointsOver}}}{LongGap.lean}{69--70}
\crow{HasDeepPoint}{\href{https://github.com/lsernau/grh-lean/blob/72fead102f663d5c7a407964a13dc0cf9d977207/grh-compact/src/New/Lift/LongGap.lean\#L75-L79}{{\leanfont Lift.\allowbreak{}KummerChart.\allowbreak{}PairingSide.\allowbreak{}HasDeepPoint}}}{LongGap.lean}{75--79}
\crow{wittFstinclKer}{\href{https://github.com/lsernau/grh-lean/blob/72fead102f663d5c7a407964a13dc0cf9d977207/grh-compact/src/New/Lift/LongGap.lean\#L88-L89}{{\leanfont Lift.\allowbreak{}KummerChart.\allowbreak{}PairingSide.\allowbreak{}wittFst\_inclKer}}}{LongGap.lean}{88--89}
\crow{isSafesum}{\href{https://github.com/lsernau/grh-lean/blob/72fead102f663d5c7a407964a13dc0cf9d977207/grh-compact/src/New/Lift/LongGap.lean\#L91-L100}{{\leanfont Lift.\allowbreak{}KummerChart.\allowbreak{}PairingSide.\allowbreak{}isSafe\_sum}}}{LongGap.lean}{91--100}
\crow{existsextendform}{\href{https://github.com/lsernau/grh-lean/blob/72fead102f663d5c7a407964a13dc0cf9d977207/grh-compact/src/New/Lift/LongGap.lean\#L105-L120}{{\leanfont Lift.\allowbreak{}KummerChart.\allowbreak{}PairingSide.\allowbreak{}exists\_extend\_form}}}{LongGap.lean}{105--120}
\crow{memofstable}{\href{https://github.com/lsernau/grh-lean/blob/72fead102f663d5c7a407964a13dc0cf9d977207/grh-compact/src/New/Lift/LongGap.lean\#L122-L124}{{\leanfont Lift.\allowbreak{}KummerChart.\allowbreak{}PairingSide.\allowbreak{}Reaches.\allowbreak{}mem\_of\_stable}}}{LongGap.lean}{122--124}
\crow{kerWithExponentsledeepElements}{\href{https://github.com/lsernau/grh-lean/blob/72fead102f663d5c7a407964a13dc0cf9d977207/grh-compact/src/New/Lift/LongGap.lean\#L126-L127}{{\leanfont Lift.\allowbreak{}KummerChart.\allowbreak{}PairingSide.\allowbreak{}kerWithExponents\_le\_deepElements}}}{LongGap.lean}{126--127}
\crow{existsinclKereqofprojTriveqzero}{\href{https://github.com/lsernau/grh-lean/blob/72fead102f663d5c7a407964a13dc0cf9d977207/grh-compact/src/New/Lift/LongGap.lean\#L132-L148}{{\leanfont Lift.\allowbreak{}KummerChart.\allowbreak{}PairingSide.\allowbreak{}exists\_inclKer\_eq\_of\_projTriv\_eq\_zero}}}{LongGap.lean}{132--148}
\crow{pointsOvernonemptyofkillsCokernel}{\href{https://github.com/lsernau/grh-lean/blob/72fead102f663d5c7a407964a13dc0cf9d977207/grh-compact/src/New/Lift/LongGap.lean\#L153-L159}{{\leanfont Lift.\allowbreak{}KummerChart.\allowbreak{}PairingSide.\allowbreak{}pointsOver\_nonempty\_of\_killsCokernel}}}{LongGap.lean}{153--159}
\crow{existsreachesdeep}{\href{https://github.com/lsernau/grh-lean/blob/72fead102f663d5c7a407964a13dc0cf9d977207/grh-compact/src/New/Lift/LongGap.lean\#L165-L191}{{\leanfont Lift.\allowbreak{}KummerChart.\allowbreak{}PairingSide.\allowbreak{}exists\_reaches\_deep}}}{LongGap.lean}{165--191}
\crow{pointsOvernonemptyofcosts}{\href{https://github.com/lsernau/grh-lean/blob/72fead102f663d5c7a407964a13dc0cf9d977207/grh-compact/src/New/Lift/LongGap.lean\#L197-L206}{{\leanfont Lift.\allowbreak{}KummerChart.\allowbreak{}PairingSide.\allowbreak{}pointsOver\_nonempty\_of\_costs}}}{LongGap.lean}{197--206}
\crow{smulsubmemkerWithExponents}{\href{https://github.com/lsernau/grh-lean/blob/72fead102f663d5c7a407964a13dc0cf9d977207/grh-compact/src/New/Lift/LongGap.lean\#L208-L211}{{\leanfont Lift.\allowbreak{}KummerChart.\allowbreak{}PairingSide.\allowbreak{}smul\_sub\_mem\_kerWithExponents}}}{LongGap.lean}{208--211}
\crow{submemkerWithExponentsofpointsOver}{\href{https://github.com/lsernau/grh-lean/blob/72fead102f663d5c7a407964a13dc0cf9d977207/grh-compact/src/New/Lift/LongGap.lean\#L213-L216}{{\leanfont Lift.\allowbreak{}KummerChart.\allowbreak{}PairingSide.\allowbreak{}sub\_mem\_kerWithExponents\_of\_pointsOver}}}{LongGap.lean}{213--216}
\crow{phipowsmulmempointsOver}{\href{https://github.com/lsernau/grh-lean/blob/72fead102f663d5c7a407964a13dc0cf9d977207/grh-compact/src/New/Lift/LongGap.lean\#L218-L227}{{\leanfont Lift.\allowbreak{}KummerChart.\allowbreak{}PairingSide.\allowbreak{}phi\_pow\_smul\_mem\_pointsOver}}}{LongGap.lean}{218--227}
\crow{wittPairingconstantAtandthirdPairing}{\href{https://github.com/lsernau/grh-lean/blob/72fead102f663d5c7a407964a13dc0cf9d977207/grh-compact/src/New/Lift/LongGap.lean\#L229-L237}{{\leanfont Lift.\allowbreak{}KummerChart.\allowbreak{}PairingSide.\allowbreak{}wittPairing\_constantAt\_and\_thirdPairing}}}{LongGap.lean}{229--237}
\crow{isSafeofexponents}{\href{https://github.com/lsernau/grh-lean/blob/72fead102f663d5c7a407964a13dc0cf9d977207/grh-compact/src/New/Lift/LongGap.lean\#L244-L259}{{\leanfont Lift.\allowbreak{}KummerChart.\allowbreak{}PairingSide.\allowbreak{}SidePairing.\allowbreak{}isSafe\_of\_exponents}}}{LongGap.lean}{244--259}
\crow{isSafesmulsub}{\href{https://github.com/lsernau/grh-lean/blob/72fead102f663d5c7a407964a13dc0cf9d977207/grh-compact/src/New/Lift/LongGap.lean\#L265-L267}{{\leanfont Lift.\allowbreak{}KummerChart.\allowbreak{}PairingSide.\allowbreak{}SidePairing.\allowbreak{}isSafe\_smul\_sub}}}{LongGap.lean}{265--267}
\crow{scaledCupBoundsofpointsatcuts}{\href{https://github.com/lsernau/grh-lean/blob/72fead102f663d5c7a407964a13dc0cf9d977207/grh-compact/src/New/Lift/LongGap.lean\#L273-L309}{{\leanfont Lift.\allowbreak{}KummerChart.\allowbreak{}PairingSide.\allowbreak{}SidePairing.\allowbreak{}scaledCupBounds\_of\_points\_at\_cuts}}}{LongGap.lean}{273--309}
\crow{pointsOvernonemptyofnoMultipleNear}{\href{https://github.com/lsernau/grh-lean/blob/72fead102f663d5c7a407964a13dc0cf9d977207/grh-compact/src/New/Lift/LongGap.lean\#L317-L323}{{\leanfont Lift.\allowbreak{}KummerChart.\allowbreak{}PairingSide.\allowbreak{}pointsOver\_nonempty\_of\_noMultipleNear}}}{LongGap.lean}{317--323}
\crow{scaledCupBoundsofkillsCokernel}{\href{https://github.com/lsernau/grh-lean/blob/72fead102f663d5c7a407964a13dc0cf9d977207/grh-compact/src/New/Lift/LongGap.lean\#L331-L350}{{\leanfont Lift.\allowbreak{}KummerChart.\allowbreak{}PairingSide.\allowbreak{}SidePairing.\allowbreak{}scaledCupBounds\_of\_killsCokernel}}}{LongGap.lean}{331--350}
\crow{raiseExponent}{\href{https://github.com/lsernau/grh-lean/blob/72fead102f663d5c7a407964a13dc0cf9d977207/grh-compact/src/New/Lift/LongGap.lean\#L359-L361}{{\leanfont Lift.\allowbreak{}KummerChart.\allowbreak{}JumpLattice.\allowbreak{}raiseExponent}}}{LongGap.lean}{359--361}
\crow{existscombofjumpMultiplicityle}{\href{https://github.com/lsernau/grh-lean/blob/72fead102f663d5c7a407964a13dc0cf9d977207/grh-compact/src/New/Lift/LongGap.lean\#L366-L404}{{\leanfont Lift.\allowbreak{}KummerChart.\allowbreak{}JumpLattice.\allowbreak{}exists\_comb\_of\_jumpMultiplicity\_le}}}{LongGap.lean}{366--404}
\crow{shiftedStd}{\href{https://github.com/lsernau/grh-lean/blob/72fead102f663d5c7a407964a13dc0cf9d977207/grh-compact/src/New/Lift/LongGap.lean\#L417-L417}{{\leanfont Lift.\allowbreak{}KummerChart.\allowbreak{}JumpLattice.\allowbreak{}shiftedStd}}}{LongGap.lean}{417--417}
\crow{constantsModE}{\href{https://github.com/lsernau/grh-lean/blob/72fead102f663d5c7a407964a13dc0cf9d977207/grh-compact/src/New/Lift/LongGap.lean\#L421-L422}{{\leanfont Lift.\allowbreak{}KummerChart.\allowbreak{}JumpLattice.\allowbreak{}constantsModE}}}{LongGap.lean}{421--422}
\crow{coeffmemconstantsModE}{\href{https://github.com/lsernau/grh-lean/blob/72fead102f663d5c7a407964a13dc0cf9d977207/grh-compact/src/New/Lift/LongGap.lean\#L430-L448}{{\leanfont Lift.\allowbreak{}KummerChart.\allowbreak{}JumpLattice.\allowbreak{}coeff\_mem\_constantsModE}}}{LongGap.lean}{430--448}
\crow{letterProjmemconstantsModE}{\href{https://github.com/lsernau/grh-lean/blob/72fead102f663d5c7a407964a13dc0cf9d977207/grh-compact/src/New/Lift/LongGap.lean\#L453-L458}{{\leanfont Lift.\allowbreak{}KummerChart.\allowbreak{}JumpLattice.\allowbreak{}letterProj\_mem\_constantsModE}}}{LongGap.lean}{453--458}
\crow{pointsOvernonemptyoflonggap}{\href{https://github.com/lsernau/grh-lean/blob/72fead102f663d5c7a407964a13dc0cf9d977207/grh-compact/src/New/Lift/LongGap.lean\#L463-L496}{{\leanfont Lift.\allowbreak{}KummerChart.\allowbreak{}PairingSide.\allowbreak{}pointsOver\_nonempty\_of\_long\_gap}}}{LongGap.lean}{463--496}}
% the notations of the defined functions: one slot for each explicit argument of the Lean declaration
\newcommand{\projY}[1]{\pr^{Y}_{#1}}
\newcommand{\tensL}[2]{\mathsf{T}_{#1,#2}}
\newcommand{\third}[1]{S^{\vee}_{#1}}
\newcommand{\pairW}[1]{\mathsf{b}_{#1}}
\newcommand{\Safe}[4]{\operatorname{Safe}_{#1}(#2;#3,#4)}
\newcommand{\CupB}[4]{\operatorname{Cup}_{#1}(#2;#3,#4)}
\newcommand{\tateL}[1]{#1(1)}
\newcommand{\inclY}[1]{\operatorname{in}^{Y}_{#1}}
\newcommand{\projCloc}[1]{\pr^{C,\mathrm{loc}}_{#1}}
\newcommand{\deepE}[3]{\operatorname{Deep}_{#1}(#2;#3)}
\newcommand{\kerX}[3]{\mathfrak{X}^{0}_{#1}(#2;#3)}
\newcommand{\ptsOver}[6]{\operatorname{Pt}_{#1}(#2,#3;#4;#5,#6)}
\newcommand{\raiseX}[5]{\operatorname{up}_{#2,#3}(#1)(#4,#5)}
\newcommand{\shStd}[4]{\operatorname{sh}_{#1}(#2)(#3,#4)}
\newcommand{\constE}[3]{\mathfrak{K}_{#1,#2,#3}}
\newcommand{\UnitMembers}{141}
\newcommand{\UnitInNote}{44}
\newcommand{\UnitProofs}{69}
\newcommand{\UnitSteps}{375}
\newcommand{\UnitAllSteps}{430}


\linespread{1.08}                       % 11 pt on 14.7 pt
\setlength{\parindent}{0pt}
\setlength{\parskip}{5.5pt plus 1.5pt minus 0.5pt}
\setlength{\abovedisplayskip}{7.5pt plus 2pt minus 3pt}
\setlength{\belowdisplayskip}{7.5pt plus 2pt minus 3pt}
\setlength{\abovedisplayshortskip}{3pt plus 2pt}
\setlength{\belowdisplayshortskip}{6pt plus 2pt minus 2pt}
\emergencystretch=1.5em
\clubpenalty=10000 \widowpenalty=10000 \displaywidowpenalty=10000
\relpenalty=7000 \binoppenalty=9000
\raggedbottom
\pagestyle{fancy}\fancyhf{}\renewcommand{\headrulewidth}{0pt}
\fancyfoot[C]{\sffamily\fontsize{7.6}{9}\selectfont\color{inkfaint}\thepage\,/\,\pageref*{LastPage}}
\setlist{topsep=2pt,partopsep=0pt,parsep=0pt,itemsep=4pt,leftmargin=2.1em,labelsep=0.55em}

\newcommand{\Q}{\mathbb{Q}}
\newcommand{\Z}{\mathbb{Z}}
\newcommand{\F}{\mathbb{F}}
\newcommand{\len}{\operatorname{len}}
\newcommand{\cok}{\operatorname{cok}}
\newcommand{\End}{\operatorname{End}}
\newcommand{\Hom}{\operatorname{Hom}}
\newcommand{\Gal}{\operatorname{Gal}}
\newcommand{\diag}{\operatorname{diag}}
\newcommand{\kw}[1]{\emph{#1}}
\newcommand{\m}[1]{\mbox{$#1$}}        % a short formula that must not be broken across lines
\newcommand{\head}[1]{\par\addvspace{17pt plus 4pt minus 2pt}\Needspace{6\baselineskip}%
  \noindent{\fontsize{13.2}{16}\selectfont\bfseries\color{ink}#1}\par\nopagebreak\vspace{-1pt}}
\newcommand{\sect}[2]{\refstepcounter{section}\label{#2}\par\addvspace{19pt plus 4pt minus 2pt}\Needspace{7\baselineskip}%
  \noindent\llap{{\fontsize{13.2}{16}\selectfont\bfseries\color{inkfaint}\thesection}\hspace{9pt}}%
  {\fontsize{13.2}{16}\selectfont\bfseries\color{ink}#1}\par\nopagebreak\vspace{-1pt}}
\newcommand{\subhead}[1]{\par\addvspace{9pt plus 2pt}\Needspace{4\baselineskip}\noindent\textbf{#1}\enspace\ignorespaces}
% where a statement lives in the Lean archive
\newcommand{\leantag}{{\sffamily\bfseries\fontsize{6.3}{8}\selectfont\color{leancol}\addfontfeatures{LetterSpace=7}LEAN}}
\newcommand{\sep}{\unskip\nobreak\hspace{5pt}{\color{inkfaint}\textperiodcentered}\hspace{5pt plus 2pt}\ignorespaces}
\newcommand{\leanline}[1]{\par\nopagebreak\vspace{-2.5pt}{\fontsize{8.9}{12.4}\selectfont\color{inksoft}\raggedright
  \noindent\llap{\leantag\hspace{9pt}}#1\par}\vspace{1pt}}
% a displayed statement in a tinted box: #1 colour, #2 text
\newcommand{\statement}[2]{\par\addvspace{7pt}\begingroup\setlength{\fboxsep}{8.5pt}\setlength{\fboxrule}{0.6pt}%
  \noindent\fcolorbox{#1!30!white}{#1!3!white}{\parbox{\dimexpr\linewidth-2\fboxsep-2\fboxrule\relax}{#2}}%
  \endgroup\par\addvspace{7pt}}
\newcommand{\nt}[2]{\par\noindent\hangindent=1.1em\hangafter=1 #1\enspace{\color{inksoft}#2}\par\vspace{1pt}}
\newcommand{\legdot}[1]{\tikz[baseline=-0.62ex]\fill[#1] (0,0) circle (2.5pt);\hspace{3pt}}
\newcommand{\crow}[4]{\csname CS@#1\endcsname & #2 & \textcolor{inksoft}{#3} & \textcolor{inksoft}{#4}\\}
\newcommand{\leanquoteline}[2]{\par\noindent\hangindent=\dimexpr#1\fontdimen2\font+2.2em\relax\hangafter=1
  \hspace*{#1\fontdimen2\font}#2}
\newcommand{\quoted}[1]{\par\vspace{3pt}{\ttfamily\fontsize{7.5}{10.4}\selectfont\color{leancol}\raggedright
  \setlength{\parskip}{0pt}\leanquote{#1}\par}\vspace{4pt}}
\newcommand{\bx}[3]{{\sffamily\bfseries\fontsize{7.3}{9.2}\selectfont\color{#1!82!black}#2}\\[1.5pt]#3}
\newcommand{\cs}[2]{\expandafter\newcommand\csname CS@#1\endcsname{#2}}

% theorems
\newtheoremstyle{plainx}{7pt plus 2pt}{7pt plus 2pt}{\itshape}{}{\bfseries}{.}{0.5em}{}
\newtheoremstyle{uprightx}{7pt plus 2pt}{7pt plus 2pt}{}{}{\bfseries}{.}{0.5em}{}
\theoremstyle{plainx}
\newtheorem{theorem}{Theorem}[section]
\newtheorem{lemma}[theorem]{Lemma}
\newtheorem{proposition}[theorem]{Proposition}
\newtheorem{corollary}[theorem]{Corollary}
\theoremstyle{uprightx}
\newtheorem{definition}[theorem]{Definition}
\newtheorem{setting}[theorem]{Setting}          % hypotheses shared by several statements; each of them cites it by number
\newtheorem{remark}[theorem]{Remark}
\newtheorem{example}[theorem]{Example}
\renewcommand{\qedsymbol}{{\setlength{\fboxsep}{0pt}\setlength{\fboxrule}{0.5pt}\fbox{\rule{0pt}{0.46em}\hspace{0.46em}}}}
% what holds throughout the note: the ambient objects only (implicit and instance arguments), never a hypothesis
\newenvironment{conventions}{\par\addvspace{5pt}\noindent\textbf{Conventions.}\enspace\ignorespaces}{\par\addvspace{3pt}}
\makeatletter
\renewenvironment{proof}[1][Proof]{\par\pushQED{\qed}\normalfont\topsep0pt\partopsep0pt
  \trivlist\item[\hskip\labelsep\itshape #1.]\ignorespaces}{\popQED\endtrivlist\@endpefalse}
\makeatother

% where each declaration of the concordance is treated: \cs{key}{place}
\cs{ofisSafe}{Theorem A}
\cs{scaledCupBoundsofpointsatcuts}{Theorem B}
\cs{scaledCupBoundsofkillsCokernel}{Theorem C}
\cs{pointsOvernonemptyoflonggap}{Theorem D}
\cs{wittFst}{Definition~\ref{d:third}}
\cs{inclKer}{Definition~\ref{d:third}}
\cs{projTrivLoc}{Definition~\ref{d:third}}
\cs{tensorLattice}{Definition~\ref{d:third}}
\cs{tateLattice}{Definition~\ref{d:third}}
\cs{thirdSide}{Definition~\ref{d:third}}
\cs{wittPairing}{Definition~\ref{d:third}}
\cs{SidePairing}{Definition~\ref{d:third}}
\cs{IsSafe}{Definition~\ref{d:safe}}
\cs{ScaledCupBounds}{Definition~\ref{d:safe}}
\cs{kerWithExponents}{Definition~\ref{d:points}}
\cs{deepElements}{Definition~\ref{d:points}}
\cs{Reaches}{Definition~\ref{d:points}}
\cs{pointsOver}{Definition~\ref{d:points}}
\cs{HasDeepPoint}{Definition~\ref{d:deep}}
\cs{raiseExponent}{Definition~\ref{d:shift}}
\cs{shiftedStd}{Definition~\ref{d:shift}}
\cs{constantsModE}{Definition~\ref{d:shift}}
\cs{wittFstsmul}{Lemma~\ref{l:comp}}
\cs{smulwittPairing}{Lemma~\ref{l:comp}}
\cs{wittFstinclKer}{Lemma~\ref{l:facts}}
\cs{memofstable}{Lemma~\ref{l:facts}}
\cs{kerWithExponentsledeepElements}{Lemma~\ref{l:facts}}
\cs{smulsubmemkerWithExponents}{Lemma~\ref{l:facts}}
\cs{submemkerWithExponentsofpointsOver}{Lemma~\ref{l:facts}}
\cs{stepBackzero}{Lemma~\ref{l:facts}}
\cs{isSafesum}{Lemma~\ref{l:sum}}
\cs{existsextendform}{Lemma~\ref{l:extend}}
\cs{existsinclKereqofprojTriveqzero}{Lemma~\ref{l:strict}}
\cs{pointsOvernonemptyofkillsCokernel}{Lemma~\ref{l:first}}
\cs{existsreachesdeep}{Proposition~\ref{p:reach}}
\cs{pointsOvernonemptyofcosts}{Lemma~\ref{l:move}}
\cs{phipowsmulmempointsOver}{Lemma~\ref{l:move}}
\cs{pointsOvernonemptyofnoMultipleNear}{Lemma~\ref{l:move}}
\cs{wittPairingconstantAtandthirdPairing}{Lemma~\ref{l:third}}
\cs{isSafeofexponents}{Proposition~\ref{p:safe}}
\cs{isSafesmulsub}{Lemma~\ref{l:safediff}}
\cs{existscombofjumpMultiplicityle}{Proposition~\ref{p:auer}}
\cs{coeffmemconstantsModE}{Lemma~\ref{l:coeff}}
\cs{letterProjmemconstantsModE}{Lemma~\ref{l:projconst}}
% ---------- common symbols of the notes on unit 136
\providecommand{\Zp}{\mathbb{Z}_p}\providecommand{\Qp}{\mathbb{Q}_p}\providecommand{\GQp}{G_{\mathbb{Q}_p}}
\providecommand{\Ainf}{A_{\inf}}\providecommand{\WC}{W(C^{\flat})}\providecommand{\BdR}{B_{\mathrm{dR}}}
\providecommand{\Fil}{\operatorname{Fil}}\providecommand{\fS}{\mathfrak{S}}\providecommand{\cE}{\mathcal{E}}
\providecommand{\cO}{\mathcal{O}}\providecommand{\fc}{\mathfrak{c}}\providecommand{\fd}{\mathfrak{d}}
\providecommand{\WW}{\mathbb{W}}\providecommand{\Per}{\operatorname{Per}}\providecommand{\pr}{\operatorname{pr}}
\providecommand{\Minf}{\mathfrak{M}^{\inf}}\providecommand{\Mloc}{\mathfrak{M}_{\mathrm{loc}}}
\providecommand{\Mlocinf}{\mathfrak{M}^{\inf}_{\mathrm{loc}}}\providecommand{\std}{\operatorname{std}}
\providecommand{\ret}{\operatorname{ret}}
\providecommand{\inputbox}[1]{\statement{inksoft}{#1}}
% ---------- notations of definitions that Part 1 owns (same forms as there)
\providecommand{\cost}[2]{\operatorname{cost}(#1,#2)}
\providecommand{\wExp}[3]{\mathfrak{X}_{#1}(#2;#3)}
\providecommand{\locLP}[2]{\mathfrak{N}_{#1,#2}}
\providecommand{\locKP}[1]{(#1)_{\fS}}
\providecommand{\invMod}[1]{\operatorname{Inv}(#1)}
\providecommand{\impE}[4]{\operatorname{imp}(#1,#2)(#3,#4)}
\providecommand{\wPair}[3]{{#1}_{W}(#2,#3)}
\providecommand{\frobX}[5]{\operatorname{fr}_{#1,#2}(#3)(#4,#5)}
\providecommand{\extL}[1]{\tilde{Y}_{#1}}
\providecommand{\trivL}[1]{C^{\mathrm{tr}}_{#1}}
\providecommand{\projC}[1]{\pr^{C}_{#1}}
\providecommand{\JumpWin}[5]{#5\subseteq[-#1,#1]\cup#2\cdot[#3,#3+#4]}
\providecommand{\raisedE}[3]{\operatorname{rs}_{#1}(#2,#3)}
\providecommand{\twoWin}[2]{\mathcal{W}_{#1}(#2)}
\providecommand{\NoMult}[3]{\operatorname{NM}(#1;#2,#3)}
\providecommand{\constAt}[3]{\operatorname{cst}_{#1,#2}(#3)}
\providecommand{\liftId}[4]{\mathfrak{a}_{#1}(#2;#3,#4)}
\providecommand{\Kills}[4]{\operatorname{Kill}_{#1}(#2;#3,#4)}
% the cost of a map (the spec says why it has no generated macro): side, cost, map
\providecommand{\Reaches}[3]{\operatorname{Cost}_{#1}(#3)\le #2}

\begin{document}
\color{ink}

{\raggedright
{\sffamily\bfseries\fontsize{7.2}{9}\selectfont\color{leancol}\addfontfeatures{LetterSpace=8}READBACK NOTE · CHAPTER VIII · UNIT 136 · PART 2 OF 3\par}
\vspace{5pt}
{\fontsize{21.5}{25.5}\selectfont\bfseries Cup products of crystalline classes vanish up to a power of~$p$: the walk along Frobenius orbits\par}
\vspace{5pt}
{\fontsize{12.4}{16}\selectfont\itshape\color{inksoft} One unit of the Lean artifact \ArtifactName, read back into
mathematics\par}}
\vspace{8pt}
{\color{rulegrey}\hrule height 0.6pt}
\vspace{7pt}
\thispagestyle{empty}

\textbf{Summary.} Two crystalline lattices with a pairing into the Tate twist modulo~$p^{M}$, and a crystalline
class of each: this note shows that the cup product of the two classes becomes a coboundary after multiplication
by an explicit power of~$p$ that does not depend on the ranks. The classes are represented by points of the
localized Kisin modules of the extensions they define; a point has exponents and a depth, and Part~1 says what it
costs to deepen it. A pair of elements is \emph{safe} if its pairing lies in the Witt vectors of the maximal ideal
of the tilt, and safety of four pairs gives the vanishing, by the Artin--Schreier--Witt sequence and a restriction
to a subgroup of index prime to~$p$. Safety is proved from exponents alone, with a point of a third extension
that carries the pairing. The points are moved along the Frobenius orbits of the letters from chosen cuts; this
is the walk. A lemma for long gaps between cuts and a relation among few elements of the same weight are proved
for use in Part~3.

\statement{leancol}{\fontsize{9.6}{13.2}\selectfont\textbf{What this note is.} The artifact \ArtifactName{} is a
body of Lean code: the compact edition of a formal proof, due to L.~Alpöge with the AI system Claude, of the Riemann
hypothesis and of its generalisation to Dirichlet $L$-functions. That headline is a claim. The authors report that
every file compiles with Lean v4.35.0 against a pinned version of Mathlib and that Lean's kernel accepted the
proofs; an independent build is recorded in the handout's record of checks when it reports. No mathematician has
read the argument through, and nothing here says that the Riemann hypothesis is proved. This note is the second of
three on one unit of the artifact, the folder on lifting; it treats \UnitInNote{} of the unit's \UnitMembers{}
declarations: the file on the walk and nine declarations of the file on the cup product bound. Theorems A to~D and
the lemmas are proved from the results of Part~1 (the note \emph{Exponents of elements of Kisin modules, the cost
of deepening them, and cokernel killers}), which are restated in a box, and from Inputs of the background library,
stated and not proved.
Every statement was compared with its Lean text, argument by argument. The Lean text read is \PinText;
Appendix~A is the concordance of \ConcDecls{} declarations.}

\sect{The statements}{s:intro}
The objects are those of Settings \ref{set:base} and~\ref{set:pair}: a Kummer chart, lattice data, sides
\m{S=(Y_S,C_S,\alpha_S)} with their extensions \m{\extL{S}}, and a pairing~$\Pi$ of two sides modulo~$p^{M}$.
Safety is Definition~\ref{d:safe}, points are Definition~\ref{d:points}, deep points are
Definition~\ref{d:deep}.

\statement{arith}{\phantomsection\label{t:A}\textbf{Theorem A} (four safe pairs give a coboundary).\enspace
\emph{In Setting~\ref{set:pair}, let $S$, $S'$ be sides over~$\cO$ and $\Pi$ a pairing of $S$ and~$S'$. Let
\m{e,e_1,e_2,e_3\ge0} be integers with \m{e_1+e_2+e_3\le e}, let \m{c\in C_S} and \m{c'\in C_{S'}}, and let
\m{x,w\in\WC\otimes\extL{S}} and \m{x'\in\WC\otimes\extL{S'}} satisfy
\[(1\otimes\projC{S})(x)=(1\otimes\projC{S})(w)=p^{e_1}(1\otimes c),\qquad
(1\otimes\projC{S'})(x')=p^{e_2}(1\otimes c').\]
Suppose that for every \m{g\in\Gamma'} the following four conditions hold:
$\Safe{\Pi}{e_3}{gx-x}{g(hx'-x')}$ for every \m{h\in\Gamma'};
\m{\Safe{\Pi}{e_3}{w-x}{gx'-x'}};
$\Safe{\Pi}{e_3}{gw-w}{g(\varphi x'-x')}$; and
$\Safe{\Pi}{e_3}{\varphi x-w}{g\varphi x'-\varphi x'}$.
Then \m{\CupB{\Pi}{e}{c}{c'}}.}}
\leanline{\lean{ofisSafe}}

\statement{arith}{\phantomsection\label{t:B}\textbf{Theorem B} (the walk).\enspace
\emph{In Setting~\ref{set:pair}, let $S$, $S'$ be sides over~$\cO$, $\Pi$ a pairing of $S$ and~$S'$,
\m{c\in C_S}, \m{c'\in C_{S'}}, and \m{\kappa,e,e'\ge0} integers. Let $\mathsf{cut}$ and $\mathsf{Q}$ be
properties of letters, and for every letter~$i$ let \m{\mathsf{l}(i)} be the least integer \m{a>0} such that
\m{i^{p^{-a}}} has the property~$\mathsf{cut}$. Let \m{o,o',z,z'} be families of exponent functions indexed by
the letters. Suppose that for every letter~$i$ with the property~$\mathsf{cut}$:
\begin{enumerate}[label=\textup{(\roman*)},leftmargin=2.3em,itemsep=2pt]
\item $S$ has a deep point over \m{p^{e}c} at~$i$ against~$S'$, in the sense of Definition~\ref{d:deep}, for
the data $(\mathsf{cut},\mathsf{l}(i),o,o',z)$ and the proposition \m{\mathsf{Q}(i)};
\item $S'$ has a deep point over \m{p^{e'}c'} at~$i$ against~$S$, for the data
$(\mathsf{cut},\mathsf{l}(i),o',o,z)$ and the proposition ``not \m{\mathsf{Q}(i)}'';
\item $\ptsOver{\third{\Pi}}{1}{i}{\kappa}{z(i)}{z'(i)}$ is not empty;
\item for every \m{\nu\ge0}, the function \m{u\mapsto z(i)(\nu,u)} is non-increasing.
\end{enumerate}
Then \m{\CupB{\Pi}{E}{c}{c'}} for every integer~$E$ with \m{e+e'+\kappa\le E}.}}
\leanline{\lean{scaledCupBoundsofpointsatcuts}}

\statement{arith}{\phantomsection\label{t:C}\textbf{Theorem C} (rank-free exponents at large weights).\enspace
\emph{In Setting~\ref{set:pair}, let $S$, $S'$ be sides over~$\cO$ and $\Pi$ a pairing of $S$ and~$S'$. Let
\m{\rho,\kappa,\eta,W,H\ge0} be integers, and for every letter~$i$ let \m{\beta(i)\ge0} be an integer and
\m{\varsigma(i)\in\{1,-1\}} a sign. Suppose:
\begin{enumerate}[label=\textup{(\roman*)},leftmargin=2.3em,itemsep=2pt]
\item for every letter~$i$: $\Kills{\third{\Pi}}{i}{\kappa}{\rho}$, \m{\Kills{S}{i}{\kappa}{\rho}} and
\m{\Kills{S'}{i}{\kappa}{\rho}};
\item for every letter~$i$: $\NoMult{p}{6(\rho+\eta+W)}{\beta(i)}$;
\item for every letter~$i$: $\JumpWin{\eta}{-\varsigma(i)}{\beta(i)}{W}{U_{Y_S}(i)}$ and
$\JumpWin{\eta}{\varsigma(i)}{\beta(i)}{W}{U_{Y_{S'}}(i)}$;
\item either no element of \m{U_{Y_S}(i)\cup U_{Y_{S'}}(i)} lies in \m{[-\rho,0)}, for any letter~$i$; or the
Kummer chart is unramified, \m{p\ge5}, \m{\eta\ge1} and $U_{Y_{S'}}(i)\subseteq[-\eta,0]$ for every
letter~$i$;
\item $3\,\bigl(\kappa+(48(\rho+\eta+W)+4)\,\fd\bigr)\le H$.
\end{enumerate}
Then \m{\CupB{\Pi}{H}{c}{c'}} for all \m{c\in C_S} and \m{c'\in C_{S'}}.}}
\leanline{\lean{scaledCupBoundsofkillsCokernel}}

\statement{arith}{\phantomsection\label{t:D}\textbf{Theorem D} (the long gap).\enspace
\emph{In Setting~\ref{set:base}, suppose that the Kummer chart is unramified. Let $S$ be a side over~$\cO$, $i$ a
letter, \m{l>0} an integer and $\kappa,\rho,\gamma,\delta,N,A\ge0$ integers such that
$\Kills{S}{i^{p^{-l}}}{\kappa}{0}$, \m{\Kills{S}{i}{0}{\rho}},
\[\begin{gathered}N=\rho+\gamma\sum_{0\le a<l}p^{a},\qquad u\le\gamma\ \text{ for all }a<l\text{ and }u\in U_{Y_S}(i^{p^{-a}}),\\
N<A,\qquad 3\,e_{\fc}\le p^{l},\qquad e_{\fc}A\le p^{l}\delta .\end{gathered}\]
Then for every \m{c\in C_S} the set
$\ptsOver{S}{c}{i}{\kappa+4\delta}{\shStd{l}{o_1}{\cdot}{\cdot}}{\shStd{l}{o_2}{\cdot}{\cdot}}$ is not empty,
where \m{o_1(u)=\min(A,-u-N)} and \m{o_2(u)=A} for \m{u\in\mathbb{Z}}.}}
\leanline{\lean{pointsOvernonemptyoflonggap}}

Theorem~C is the result that Part~3 uses most: the exponent~$H$ depends on the windows and on the killers, not
on the ranks of the lattices. Theorem~D produces a deep point when the cuts are far apart.

\sect{The setting}{s:setting}
\begin{conventions}\phantomsection\label{c:conventions}
Throughout, $p$ is a prime number. The declarations of this unit take the following ambient objects as implicit or
instance arguments; they are fixed by Settings \ref{set:base} and~\ref{set:pair} and named again in each
statement that needs them: the prime~$p$; the group~$\Gamma$ over~$\GQp$; the Kummer chart~$\fc$ for~$\Gamma$; the
coefficient ring~$\cO$; a lattice datum~$L$; sides $S$, $S'$; a pairing~$\Pi$; elements \m{c\in C_S} and
\m{c'\in C_{S'}}; a letter~$i$; integers $l$, $\kappa$, $\rho$, $e$, $e'$, $n$, $A$; exponent functions
$o$, $o'$, $o_1$, $o_2$, $o_3$; elements $x$, $w$, $y$ of the modules \m{\WC\otimes\extL{S}}; an element~$g$
of~$\Gamma'$. Tensor products without a subscript are over~$\Zp$; \m{v_p(n)} is the exponent of~$p$ in an
integer~$n$, with \m{v_p(0)=0}. An \kw{exponent function} is a function
$o\colon\mathbb{N}\times\mathbb{Z}\to\mathbb{Z}$; inequalities and minima of exponent functions are pointwise.
The Lean text calls the artifact's period ring \emph{Bcris}; it is Colmez's ring~$B_{\max}$.
\end{conventions}

\begin{setting}[the base: chart, lattice data, sides; from the background library and from Part~1]\label{set:base}
The notation is that of Settings 2.1 and~2.2 and Definitions 2.3 to~2.8 of Part~1, which this Setting recalls.

\emph{Rings and the chart.} \m{\Ainf\subset\WC} are the Witt vectors of the tilt~$\mathcal{O}_C^\flat$ of the
integers of the completed algebraic closure~$C$ of~$\Qp$ and of its fraction field, with Frobenius~$\varphi$;
\m{W(\mathfrak{m}^\flat)} is the ideal of~$\Ainf$ of the Witt vectors with components in the maximal
ideal~$\mathfrak{m}^\flat$ of the tilt; $\theta\colon\Ainf\to\mathcal{O}_C$ is Fontaine's map; $\BdR$ is the
field of de Rham periods with its filtration \m{\Fil^{n}\BdR}, \m{n\in\mathbb{Z}}; \m{\tau=[\varepsilon^{1/p}]-1}.
$\Gamma$ is a group with a homomorphism \m{j_\Gamma\colon\Gamma\to\GQp} of open image, and $\chi$ its cyclotomic
character; \m{\mathbb{Z}\times\Gamma} acts on~$\WC$ through~$\varphi$ and~$j_\Gamma$. A Kummer chart~$\fc$
for~$\Gamma$ is fixed: its ring \m{\fS\subset\Ainf}, a quotient of \m{\WW[[X]]} with \m{\WW=W(\mathbb{F}_{p^h})},
its coordinate~$z$, the Teichmüller lift of a non-unit of the tilt, its Eisenstein element~$\cE$ with
\m{\cE_\nu=\varphi^{\nu}(\cE)}, its index~$e_{\fc}$ and the number~$\fd$, its subgroup~$\Gamma'$ of index prime
to~$p$, which contains all elements near~$1$, the field~$F$ fixed by~$\Gamma'$, the fields
\m{F_\nu\subset\BdR}, and the retraction $\operatorname{co}_0\colon\Ainf\to\fS$. The chart is \kw{unramified} if
\m{\cE=z+p}, \m{e_{\fc}=1}, \m{\fd=0}. The \kw{letters} are the elements of~\m{\mathbb{F}_{p^h}^\times}, with
Teichmüller lifts~$[i]$.

\emph{Lattice data.} $\cO$ is the ring of integers of a finite extension of~$\Qp$ with a root of
unity~$\omega$, \m{\omega^{p^h-1}=1}. A lattice datum~$L$ over~$\cO$ is a finite free $\cO$-module with an
$\cO$-linear action of~$\Gamma$, crystalline on~$\Gamma'$, with finite sets \m{U_L(i)\subset\mathbb{Z}} that
contain its Hodge--Tate jumps at the embeddings above the letter~$i$. It has its module \m{\WC\otimes L} with
the action of \m{\mathbb{Z}\times\Gamma}, the lattice \m{\Ainf\otimes L}, the parts \m{(\WC\otimes L)_i} with
projectors~$\pr_i$, the periods \m{\Per(L)} and \m{\Per_{i,u}(L)}, the Kisin modules \m{\Minf(L)},
\m{\Mlocinf(L)}, \m{\Mloc(L)}, the retraction $\ret_L\colon\Mlocinf(L)\to\Mloc(L)$, and for a letter~$i$ and
an exponent function~$o$ the module \m{\wExp{L}{i}{o}} of the elements of
$\Mlocinf(L)\cap(\WC\otimes L)_i$ with
$\mu(\varphi^{-\nu}x)\in\Fil^{o(\nu,u)}\BdR$ for all \m{\nu\ge0}, \m{u\in U_L(i^{p^{-\nu}})} and
\m{\mu\in\Per_{i^{p^{-\nu}},u}(L)}. Further $\locLP{l}{i}(L)=\Mloc(L)\cap\wExp{L}{i}{\std_l}$, with
\m{\std_l(\nu,u)=0} for \m{\nu<l} and \m{-u} for \m{\nu\ge l}; \m{\locKP{Q}} is the set of the elements of
\m{\Mloc(L)} in a submodule~$Q$; \m{\invMod{Q}} is the set of the~$x$ with \m{gx-x\in Q} for all
\m{g\in\Gamma'}; \m{\cost{\nu}{n}=v_p(n)+\fd+3\nu}; $\impE{o^{\mathrm{lo}}}{o^{\mathrm{hi}}}{\nu}{u}$ is
\m{o^{\mathrm{hi}}(\nu,u)} if \m{o^{\mathrm{lo}}(\nu,u)>0} and \m{\min(o^{\mathrm{hi}}(\nu,u),0)} otherwise;
\m{\frobX{l}{n}{o}{\nu}{u}} is~$n$ for \m{\nu<l} and \m{o(\nu-l,u)} for \m{\nu\ge l}; and
\m{\wPair{b}{y}{x}} is the value of the $\WC$-bilinear extension of an $\cO$-bilinear map~$b$ of lattices.

\emph{Sides.} A side \m{S=(Y_S,C_S,\alpha_S)} over~$\cO$ consists of a lattice datum~$Y_S$ with
\m{0\in U_{Y_S}(i)} for all letters~$i$, a lattice~$C_S$, and an $\cO$-linear map~$\alpha_S$ from~$C_S$ to the
crystalline cocycles \m{Z^{1}_{\mathrm{f}}(Y_S)}. Its extension \m{\extL{S}} is \m{Y_S\times C_S} with
\m{g(y,c)=(gy+\alpha_S(c)(g),c)} and the sets \m{U_{Y_S}(i)}; \m{\trivL{S}} is $C_S$ with the trivial action and
the sets~$\{0\}$; $\projC{S}\colon\extL{S}\to\trivL{S}$ is the projection. For \m{c\in C_S},
$\constAt{S}{i}{c}=\pr_i(1\otimes c)$; \m{\liftId{S}{r}{l}{i}} is the lifting ideal, and
\m{\Kills{S}{i}{\kappa}{\rho}} means $p^{\kappa}\cE^{\rho}\in\liftId{S}{0}{0}{i}$. Finally
\m{\raisedE{n}{\nu}{u}} is \m{n-u} for \m{\nu=0} and \m{-u} for \m{\nu>0};
$\twoWin{n}{c}=[-n,n]\cup[c-n,c+n]$; \m{\NoMult{p}{n}{\beta}} says that $p$ divides no non-zero element of
$\twoWin{n}{\beta}\cup\twoWin{n}{-\beta}$; and $\JumpWin{\eta}{\varsigma}{\beta}{W}{U}$ says that every
\m{u\in U} has \m{|u|\le\eta} or \m{\beta\le\varsigma u\le\beta+W}.
\end{setting}

\begin{setting}[pairings and the cup product; from the background library]\label{set:pair}
In Setting~\ref{set:base}:

\emph{Twists and duals.} For a lattice~$T$ over~$\cO$, $T^{\vee}(1)=\mathrm{Hom}_{\cO}(T,\cO(1))$ is the twisted
dual, where \m{\cO(1)} is $\cO$ with $\Gamma$ acting through~$\chi$; for lattices $T$, $T'$,
\m{T\otimes_{\cO}T'} carries the diagonal action.

\emph{Pairings.} For lattices $Y$, $Y'$ over~$\cO$, a \kw{pairing} of $Y$ and~$Y'$ consists of an integer
\m{M\ge0}, an $\cO$-linear map $\psi\colon Y\otimes_{\cO}Y'\to\cO(1)$, a $\Zp$-linear map
$\operatorname{tr}\colon\cO\to\Zp$, and a crystalline cocycle
$c_\Pi\in Z^{1}_{\mathrm{f}}((Y\otimes_{\cO}Y')^{\vee}(1))$ with \m{p^{M}c_\Pi(g)=g\psi-\psi} for all
\m{g\in\Gamma}. Its \kw{pairing modulo~$p^{M}$} is
\[\langle y,y'\rangle=\operatorname{tr}(\psi(y\otimes y'))\bmod p^{M}\ \in\ \mathbb{Z}/p^{M},\qquad
\langle gy,gy'\rangle=\chi(g)\,\langle y,y'\rangle .\]
A \kw{pairing of sides} $S$ and~$S'$ is a pairing of the lattices $Y_S$ and~$Y_{S'}$. The
\kw{third pairing} of~$\Pi$ is the $\cO$-bilinear map
\[\mathsf{p}_\Pi\colon\bigl((Y\otimes_{\cO}Y')^{\vee}(1)\times\cO\bigr)\times(Y\otimes_{\cO}Y')\to\cO(1),\qquad
((\lambda,a),t)\mapsto p^{M}\lambda(t)+a\,\psi(t);\]
it is $\Gamma$-equivariant when the first factor carries the action of the extension of~$\cO$ along
\m{a\mapsto a\,c_\Pi}.

\emph{The module of values.} $\mathbb{M}_M=\WC\otimes\mathbb{Z}/p^{M}(1)$, with \m{\mathbb{Z}\times\Gamma}
acting on the second factor through~$\chi$; its \kw{small part}~$\mathfrak{s}_M$ is the image of
$W(\mathfrak{m}^\flat)\cdot(1\otimes1)$, an $\Ainf$-submodule stable under \m{\mathbb{Z}\times\Gamma}.

\emph{The cup product.} Fix an injective homomorphism~$\mathsf{e}_M$ from \m{\mathbb{Z}/p^{M}} to the
multiplicative group of~$\bar{\mathbb{Q}}_p$. For a topology on~$\Gamma$ and maps \m{a\colon\Gamma\to Y},
\m{a'\colon\Gamma\to Y'}, the \kw{cup product of $a$ and~$a'$ under~$\Pi$ bounds} if the function
$(g,h)\mapsto\mathsf{e}_M\bigl(\langle a(g),g\,a'(h)\rangle\bigr)$ on \m{\Gamma\times\Gamma} is the coboundary of a
continuous $1$-cochain of~$\Gamma$ with values in~$\bar{\mathbb{Q}}_p^{\times}$, on which $\Gamma$ acts
through~$j_\Gamma$.
\end{setting}

\begin{definition}[the objects attached to a pairing of sides]\label{d:third}
In Setting~\ref{set:pair}, let $S$, $S'$ be sides over~$\cO$ and $\Pi$ a pairing of $S$ and~$S'$, modulo~$p^{M}$.
\begin{enumerate}[label=\textup{(\roman*)},leftmargin=2.3em,itemsep=2pt]
\item $\projY{S}\colon\WC\otimes\extL{S}\to\WC\otimes Y_S$ is the $\WC$-linear map induced by
\m{(y,c)\mapsto y}, and \m{\inclY{S}\colon Y_S\to\extL{S}} is the equivariant map \m{y\mapsto(y,0)}.
\item \m{\projCloc{S}} is the $\Zp$-linear map $\Mloc(\extL{S})\to\WC\otimes C_S$,
\m{x\mapsto(1\otimes\projC{S})(x)}.
\item \m{\tensL{S}{S'}} is the lattice datum \m{Y_S\otimes_{\cO}Y_{S'}} with the sets
\m{U_{Y_S}(i)+U_{Y_{S'}}(i)} of sums.
\item \m{\tateL{\cO}} is the lattice datum on~$\cO$ with $\Gamma$ acting through~$\chi$ and the set~$\{-1\}$ at
every letter.
\item \m{\third{\Pi}} is the side whose lattice datum is $(Y_S\otimes_{\cO}Y_{S'})^{\vee}(1)$ with the sets
$\{0\}\cup\{-v-1: v\in U_{Y_S}(i)+U_{Y_{S'}}(i)\}$, whose lattice is~$\cO$ with the trivial action, and
whose cocycle map is \m{a\mapsto a\,c_\Pi}.
\item \m{\pairW{\Pi}} is the $\WC$-bilinear extension
$(\WC\otimes Y_S)\times(\WC\otimes Y_{S'})\to\mathbb{M}_M$ of the pairing modulo~$p^{M}$; its value at
\m{(y,y')} is written \m{\pairW{\Pi}(y,y')}.
\end{enumerate}
That these are lattice data and a side rests on Input~\ref{j:class}. The type of the pairings of $S$
and~$S'$ is a declaration of its own in the formal text.
\end{definition}
\leanline{\lean{wittFst}\sep \lean{inclKer}\sep \lean{projTrivLoc}\sep \lean{tensorLattice}\sep \lean{tateLattice}\sep
\lean{thirdSide}\sep \lean{wittPairing}\sep \lean{SidePairing}}

\begin{definition}[safe pairs; cup products that bound after scaling]\label{d:safe}
In Setting~\ref{set:pair}, let $S$, $S'$ be sides over~$\cO$ and $\Pi$ a pairing of $S$ and~$S'$,
modulo~$p^{M}$, and let \m{e\ge0} be an integer.
\begin{enumerate}[label=\textup{(\roman*)},leftmargin=2.3em,itemsep=2pt]
\item For \m{w\in\WC\otimes\extL{S}} and \m{y\in\WC\otimes\extL{S'}}, \m{\Safe{\Pi}{e}{w}{y}} means
\[p^{e}\,\pairW{\Pi}\bigl(\projY{S}(w),\projY{S'}(y)\bigr)\in\mathfrak{s}_M .\]
\item For \m{c\in C_S} and \m{c'\in C_{S'}}, \m{\CupB{\Pi}{e}{c}{c'}} means that, for every topology on~$\Gamma$
for which $j_\Gamma$ is an open embedding, the cup product of \m{\alpha_S(c)} and \m{p^{e}\alpha_{S'}(c')}
under~$\Pi$ bounds.
\end{enumerate}
\end{definition}
\leanline{\lean{IsSafe}\sep \lean{ScaledCupBounds}}

\begin{definition}[deep elements, costs, points]\label{d:points}
In Setting~\ref{set:base}, let $S$ be a side over~$\cO$, $i$ a letter and $o$, $o_1$, $o'$ exponent functions.
\begin{enumerate}[label=\textup{(\roman*)},leftmargin=2.3em,itemsep=2pt]
\item \m{\kerX{S}{i}{o}} is the $\Ainf$-module of the \m{x\in\wExp{\extL{S}}{i}{o}} with
\m{(1\otimes\projC{S})(x)=0}.
\item \m{\deepE{S}{i}{o}} is the $\Zp$-module of the \m{x\in\Mloc(\extL{S})} such that
$(1\otimes\projC{S})(x)=\constAt{S}{i}{c}$ for some \m{c\in C_S} and
$x\in\invMod{\wExp{\extL{S}}{i}{o}}$.
\item For an integer \m{e\ge0} and an additive map $\mathsf{A}\colon\Mloc(\extL{S})\to\Mloc(\extL{S})$,
\m{\Reaches{S}{e}{\mathsf{A}}} means that
$\mathsf{A}(x)-p^{e}x\in\locKP{\kerX{S}{j}{o''}}$ for every letter~$j$, every exponent function~$o''$ and
every \m{x\in\deepE{S}{j}{o''}}; in words, $\mathsf{A}$ \kw{has cost~$e$}.
\item For \m{c\in C_S} and an integer \m{e\ge0}, \m{\ptsOver{S}{c}{i}{e}{o_1}{o'}} is the set of the
\m{x\in\Mloc(\extL{S})} with $(1\otimes\projC{S})(x)=p^{e}\,\constAt{S}{i}{c}$,
\m{x\in\wExp{\extL{S}}{i}{o_1}} and $x\in\invMod{\wExp{\extL{S}}{i}{o'}}$: the \kw{points over~$p^{e}c$
at~$i$ with exponents~$o_1$ and depth~$o'$}.
\end{enumerate}
\end{definition}
\leanline{\lean{kerWithExponents}\sep \lean{deepElements}\sep \lean{Reaches}\sep \lean{pointsOver}}

\begin{definition}[a deep point against another side]\label{d:deep}
In Setting~\ref{set:base}, let $S$, $S'$ be sides over~$\cO$, \m{c\in C_S}, $i$ a letter, \m{e,l\ge0} integers,
$\mathsf{cut}$ a property of letters, \m{o,o',z} families of exponent functions indexed by the letters, and
$\mathsf{q}$ a proposition. \emph{$S$ has a deep point over~$p^{e}c$ at~$i$ against~$S'$, for the data
\m{(\mathsf{cut},l,o,o',z)} and the proposition~$\mathsf{q}$}, if there are an exponent function~$d$ and an
integer \m{r\ge0} such that \m{\ptsOver{S}{c}{i}{e}{o(i)}{d}} is not empty, \m{i^{p^{-r}}} has the
property~$\mathsf{cut}$, and for all \m{\nu\ge0}, \m{u\in U_{Y_S}(i^{p^{-\nu}})} and
\m{v\in U_{Y_{S'}}(i^{p^{-\nu}})}:
\begin{enumerate}[label=\textup{(\alph*)},leftmargin=2.3em,itemsep=2pt]
\item $-z(i)(\nu,-u-v-1)-d(\nu,u)<\min\bigl(-v,\ o'(i)(\nu,v)\bigr)$;
\item if $\mathsf{q}$ holds, then \m{\std_l(\nu,u)\le o(i)(\nu,u)} and
\[-\frobX{r}{0}{z(i^{p^{-r}})}{\nu}{-u-v-1}-d(\nu,u)<
\min\bigl(o'(i)(\nu,v),\ \frobX{l}{0}{o'(i^{p^{-l}})}{\nu}{v}\bigr).\]
\end{enumerate}
\end{definition}
\leanline{\lean{HasDeepPoint}}

\begin{definition}[exponents for the long gap]\label{d:shift}
In Setting~\ref{set:base}:
\begin{enumerate}[label=\textup{(\roman*)},leftmargin=2.3em,itemsep=2pt]
\item For an exponent function~$o$, integers \m{\nu,\nu'\ge0} and \m{u,v\in\mathbb{Z}},
$\raiseX{o}{\nu}{u}{\nu'}{v}=o(\nu',v)+1$ if \m{(\nu',v)=(\nu,u)}, and
$\raiseX{o}{\nu}{u}{\nu'}{v}=o(\nu',v)$ otherwise.
\item For an integer \m{l\ge0}, a function $o_1\colon\mathbb{Z}\to\mathbb{Z}$, \m{\nu\ge0} and
\m{u\in\mathbb{Z}}, $\shStd{l}{o_1}{\nu}{u}=\std_l(\nu,u)+o_1(u)$ if \m{\nu=0}, and
$\shStd{l}{o_1}{\nu}{u}=\std_l(\nu,u)$ otherwise.
\item For a lattice datum~$L$, integers \m{l,A\ge0} and a letter~$i$, \m{\constE{l}{A}{i}(L)} is the
$\Zp$-module of the \m{x\in\locLP{l}{i}(L)} with $x\in\invMod{\wExp{L}{i}{\shStd{l}{o_2}{\cdot}{\cdot}}}$,
where \m{o_2} is the constant function with value~$A$.
\end{enumerate}
\end{definition}
\leanline{\lean{raiseExponent}\sep \lean{shiftedStd}\sep \lean{constantsModE}}

\inputbox{\textbf{Inputs from Part~1.} The following results of the note \emph{Exponents of elements of Kisin
modules, the cost of deepening them, and cokernel killers} are used; $L$ is a lattice datum, $S$ a side, $i$ a
letter, $o$, $o'$ exponent functions, \m{g\in\Gamma'}.
\begin{enumerate}[label=\textup{P\arabic*.},ref=P\arabic*,leftmargin=2.7em,itemsep=2.5pt]
\item\label{q:loc} \emph{Localized modules} (Lemma~4.1). \m{\cE_\nu^{\,n}x\in\Mloc(L)} if and only if
\m{x\in\Mloc(L)}; \m{\Mloc(L)} is stable under~$\varphi$; \m{\Mlocinf(L)} is stable under~$\Gamma'$ and is the
$\Ainf$-span of \m{\Mloc(L)}; \m{1\otimes f} maps \m{\Mloc(L)} into \m{\Mloc(L')} for an equivariant~$f$; and
$\mu(\varphi^{-\nu}(gx-x))=g\mu(\varphi^{-\nu}x)-\mu(\varphi^{-\nu}x)$ for \m{x\in\Mlocinf(L)},
\m{\mu\in\Per(L)}.
\item\label{q:exp} \emph{Exponent modules} (Lemma~4.2). $\wExp{L}{i}{o}\subseteq\wExp{L}{i}{o'}$ if
\m{o'\le o} on the jump sets; to pass from $o$ to a larger~$o'$ it suffices to check the values where
\m{o<o'}; \m{\cE_{\nu_0}^{\,n}} maps \m{\wExp{L}{i}{o}} into the module with exponents raised by~$n$
at~$\nu_0$; \m{gx-x\in\wExp{L}{i}{o}} for \m{x\in\wExp{L}{i}{o}}; and an equivariant \m{f\colon L\to L'} with
\m{U_{L'}(j)\subseteq U_L(j)} for all~$j$ maps \m{\wExp{L}{i}{o}} into \m{\wExp{L'}{i}{o}}.
\item\label{q:ztau} \emph{Division by~$z\tau$} (Lemma~4.3). If \m{x\in\Mloc(L)} and
\m{gx-x\in\wExp{L}{i}{o}}, then $(z\tau)^{-1}(gx-x)\in\wExp{L}{i}{o}$.
\item\label{q:any} \emph{Values at any type} (Lemma~4.4). If \m{x\in\wExp{L}{i}{o}},
\m{\mu\in\Per_{i^{p^{-\nu}},u}(L)} and a $B^{+}_{\mathrm{dR}}$-module~$M$ contains \m{\Fil^{o(\nu,v)}\BdR} for
all \m{v\in U_L(i^{p^{-\nu}})} with \m{v\le u}, then \m{\mu(\varphi^{-\nu}x)\in M}.
\item\label{q:ret} \emph{The retraction} (Lemmas 5.1 and~5.3). \m{\ret_L} is $\fS$-linear, the identity on
\m{\Mloc(L)}, with $\ret_L(ax)=\operatorname{co}_0(a)x$ for \m{a\in\Ainf}, \m{x\in\Mloc(L)}; it commutes with
\m{1\otimes f} for equivariant~$f$; it maps \m{\wExp{L}{i}{o}} into itself; and if
$\mu(\varphi^{-l}y)-v\in\Fil^{n}\BdR$ with \m{v\in F_l}, then
$\mu(\varphi^{-l}\ret_L(y))-v\in\Fil^{n}\BdR$.
\end{enumerate}}

\inputbox{\textbf{Inputs from Part~1, continued.}
\begin{enumerate}[label=\textup{P\arabic*.},ref=P\arabic*,leftmargin=2.7em,itemsep=2.5pt,start=6]
\item\label{q:deep} \emph{Deepening} (Lemma~6.2 and Proposition~6.3). If
$x\in\wExp{L}{i}{o^{\mathrm{lo}}}\cap\invMod{\wExp{L}{i}{o^{\mathrm{hi}}}}$, then
$x\in\wExp{L}{i}{\impE{o^{\mathrm{lo}}}{o^{\mathrm{hi}}}{\cdot}{\cdot}}$. For integers \m{\nu\ge0} and~$n$
there are \m{t_1,\dots,t_m\in\Gamma'} and \m{y\in\Ainf} such that, for \m{x\in\Mloc(L)} and
\m{\mu\in\Per(L)}: if $\mu(\varphi^{-\nu}x)-c\in\Fil^{n}\BdR$ for some \m{c\in F}, then
$\mu(\varphi^{-\nu}(p^{\cost{\nu}{n}}x+\sum_ky(z\tau)^{-1}(t_kx-x)))-c'\in\Fil^{n+1}\BdR$ for some
\m{c'\in F}.
\item\label{q:B} \emph{Theorem~B and its corollaries} (Theorem~B, Lemmas 7.2 and~7.3).
\m{x\in\wExp{L}{i}{\std_l}} if and only if \m{\varphi^{-l}x\in\Minf(L)} and \m{x\in(\WC\otimes L)_i}. If
\m{x\in\wExp{L}{i^{p^{-l}}}{o}}, then $\varphi^{l}x\in\wExp{L}{i}{\frobX{l}{0}{o}{\cdot}{\cdot}}$. If $x$ is a
point over~$F_0$ of \m{\varphi^{l}\Minf(L)}, then \m{\pr_i(x)\in\locLP{l}{i}(L)}.
\item\label{q:C} \emph{Theorem~C} (exponents add). For an equivariant $\cO$-bilinear map
\m{b\colon A\times L\to C} of lattice data, \m{y\in\wExp{A}{i}{o_A}} and \m{x\in\wExp{L}{i}{o}}: if
$o_C(\nu,v)\le o(\nu,u')+o_A(\nu,u)$ whenever \m{u\in U_A(i^{p^{-\nu}})}, \m{u'\in U_L(i^{p^{-\nu}})},
\m{v\in U_C(i^{p^{-\nu}})} and \m{u+u'\le v}, then $\wPair{b}{y}{x}\in\wExp{C}{i}{o_C}$. Moreover
$\wPair{b}{\pr_i(y)}{x}=\pr_i(\wPair{b}{y}{x})$.
\item\label{q:kill} \emph{The lifting ideal} (Lemmas 9.2, 9.3 and~9.4, Theorem~D). If $\Gamma$ acts on~$L$
through~$\chi^{c}$ and \m{x\in\wExp{L}{i}{\std_0}}, then \m{\tau^{c}x\in\Ainf\otimes L}.
$\constAt{S}{i}{c}\in\locLP{l}{i}(\trivL{S})$ and $\varphi^{l}\constAt{S}{i}{c}=\constAt{S}{i^{p^{l}}}{c}$.
$\varphi^{l}(s)\in\liftId{S}{r}{l}{i^{p^{l}}}$ if and only if \m{s\in\liftId{S}{r}{0}{i}}. If \m{f\in\fS}
lies in \m{\liftId{S}{0}{l}{i}}, then every \m{y\in\locLP{l}{i}(\trivL{S})} has an
\m{x\in\locLP{l}{i}(\extL{S})} with \m{(1\otimes\projC{S})(x)=fy}. If the chart is unramified, \m{p\ge5} and
all \m{u\in U_{Y_S}(j)} are \m{\le(p-3)(\kappa+\rho)}, then \m{\Kills{S}{i}{\kappa}{\rho}}.
\end{enumerate}}

\inputbox{\textbf{Inputs from the background library.} Each is a declaration of the background library, read in its
Lean text; the number in brackets is the unit that owns it.
\begin{enumerate}[label=\textup{J\arabic*.},ref=J\arabic*,leftmargin=2.7em,itemsep=2.5pt]
\item\label{j:class} \emph{Stability of the class} [81, 86]. The tensor product of two lattice data, the twisted
dual of a lattice datum with the negatives of its jump sets shifted by~$-1$, and \m{\cO(c)} with the
set~$\{-c\}$ are lattice data. The third pairing~$\mathsf{p}_\Pi$ is $\Gamma$-equivariant. Reference to be
confirmed.
\item\label{j:asw} \emph{Artin--Schreier--Witt over the tilt} [82]. Let \m{M\ge0}. \emph{(a)} $p^{M}$
kills~$\mathbb{M}_M$, and every \m{d\in\mathbb{M}_M} is \m{\varphi d'-d'} for some~$d'$. \emph{(b)} If
\m{\varphi d-d\in\mathfrak{s}_M}, then \m{d-1\otimes k\in\mathfrak{s}_M} for some
\m{k\in\mathbb{Z}/p^{M}}. \emph{(c)} If \m{1\otimes k\in\mathfrak{s}_M}, then \m{k=0}. \emph{(d)} Let $T$ be
a $\Zp$-representation of~$\Gamma$ such that for each \m{t\in T}, \m{gt-t\in p^{M}T} for $g$ near~$1$, and
let $\Lambda\colon\WC\otimes T\to\mathbb{M}_M$ be $\WC$-linear. Then for each \m{v\in\WC\otimes T},
\m{\Lambda(gv-v)\in\mathfrak{s}_M} for $g$ near~$1$. In the handout's catalogue this is a weaker form of the
first step of Herr's complex (L.~Herr, \emph{Sur la cohomologie galoisienne des corps $p$-adiques}, Bull.\ Soc.\
Math.\ France~126 (1998), 563--600; place not confirmed); in the second reader's comparison the
Artin--Schreier--Witt sequence (L.~Illusie, \emph{Complexe de de Rham--Witt}, §2.7) explains~(a), and the
statements about the small part and~(d) are not located in the literature.
\item\label{j:cont} \emph{Crystalline cocycles are small near~$1$} [87]. For a potentially crystalline
lattice~$T$, \m{c\in Z^{1}_{\mathrm{f}}(T)} and \m{M\ge0}: \m{c(g)\in p^{M}T} for $g$ near~$1$. For
\m{t\in T} the map \m{g\mapsto gt-t} lies in \m{Z^{1}_{\mathrm{f}}(T)}. Reference to be confirmed.
\item\label{j:res} \emph{Restriction and the index} [95]. For an open subgroup~$\Delta$ of finite index in a
topological group and a class in the second continuous cohomology whose restriction to~$\Delta$ vanishes, the
index times the class is~$0$. (Corestriction after restriction is multiplication by the index: J.-P.~Serre,
\emph{Local Fields}, Ch.~VII.)
\end{enumerate}}

\inputbox{\textbf{Inputs from the background library, continued.}
\begin{enumerate}[label=\textup{J\arabic*.},ref=J\arabic*,leftmargin=2.7em,itemsep=2.5pt,start=5]
\item\label{j:tate} \emph{Tate} [58]. If \m{V\in\BdR}, $M$ is a $B^{+}_{\mathrm{dR}}$-submodule of~$\BdR$ and
\m{gV-V\in M} for all \m{g\in\Gamma'}, there is \m{c\in F} with \m{V-c\in M}. (A weaker form of Tate's
theorems; see Part~1, Input~I6.)
\item\label{j:types} \emph{Periods of one type} [82]. For a letter~$i$ and \m{u\in\mathbb{Z}} there are
\m{r\le m_L(i,u)} periods $\mu_1,\dots,\mu_r\in\Per_{i,u}(L)$ such that every \m{\mu\in\Per_{i,u}(L)} is a
combination of them with coefficients in~$\bar{\mathbb{Q}}_p$, modulo \m{\Per_{i,u-1}(L)}; here \m{m_L(i,u)} is the total
multiplicity of the jump~$u$ of~$L$ over the embeddings above~$i$. The handout's catalogue calls it the authors'
analogue for Kisin modules of statements on Wach modules (L.~Berger, \emph{Limites de représentations
cristallines}, Compositio Math.~140 (2004), §III; N.~Wach, Bull.\ Soc.\ Math.\ France~124 (1996); place not
confirmed). Also: every $\Zp$-linear map \m{L\to\BdR} extends to an $\Ainf$-linear map on \m{\WC\otimes L}, and
\m{\mu\circ\pr_i\in\Per(L)} for \m{\mu\in\Per(L)}.
\item\label{j:auer} \emph{Auerbach's lemma} [general library]. Let $\mathsf{U}$ be a finite-dimensional space
of functions with values in~$C$ on a finite set. Then there are points \m{s_1,\dots,s_q} of the set, $q$ the
dimension, and a basis $\mathsf{z}_1,\dots,\mathsf{z}_q$ of~$\mathsf{U}$ such that
\m{\mathsf{z}_b(s_{b'})} is $1$ for \m{b=b'} and $0$ otherwise, and \m{|\mathsf{z}_b(k)|\le1} for every
point~$k$; so $\mathsf{v}(k)=\sum_b\mathsf{z}_b(k)\,\mathsf{v}(s_b)$ for all \m{\mathsf{v}\in\mathsf{U}}. This is
a special case of the declaration, which allows a compact set of linear forms. (H.~Auerbach, thesis, Lwów
1930; place not confirmed; in the second reader's comparison a determinant-maximization extension of Auerbach's
lemma.)
\item\label{j:frame} \emph{The frame theorem: spanning} [83]. Let \m{l,A,\delta\ge0} with
\m{3e_{\fc}\le p^{l}} and \m{e_{\fc}A\le p^{l}\delta}, and let $x$ be a point over~$F_0$ of
\m{\varphi^{l}\Minf(L)}. Then \m{p^{4\delta}x} is a combination, with coefficients in~$\fS$, of points~$y$
over~$F_0$ of \m{\varphi^{l}\Minf(L)} with \m{g\mu(y)-\mu(y)\in\Fil^{A}\BdR} for all \m{\mu\in\Per(L)},
\m{g\in\Gamma'}. Reference to be confirmed (the file calls the frame theorem a known result: the Kisin lattice
of a crystalline lattice).
\item\label{j:mod} \emph{Residues} [83]. If the chart is unramified, every \m{s\in\fS} is \m{w+\cE t} with
\m{w\in\WW} and \m{t\in\fS}. The ring~$\WW$ is free of rank~$h$ over~$\Zp$, its elements are fixed
by~$\Gamma'$, and the difference of the Teichmüller lifts of two distinct letters is a unit. The
coordinate~$z$ lies in \m{W(\mathfrak{m}^\flat)}.
\end{enumerate}}

\sect{Four safe pairs give a coboundary}{s:safe}
Throughout this section the setting is Setting~\ref{set:pair}; $S$, $S'$ are sides over~$\cO$ and $\Pi$ is a
pairing of $S$ and~$S'$, modulo~$p^{M}$.

\begin{lemma}[the action on the components]\label{l:comp}
Let $n=(a,g)\in\mathbb{Z}\times\Gamma$.
\begin{enumerate}[label=\textup{(\alph*)},leftmargin=2.3em,itemsep=2pt]
\item For \m{x\in\WC\otimes\extL{S}},
$\projY{S}(n\cdot x)=n\cdot\projY{S}(x)+(1\otimes\alpha_S(\cdot)(g))\bigl(n\cdot(1\otimes\projC{S})(x)\bigr)$,
where $\alpha_S(\cdot)(g)\colon C_S\to Y_S$ is the map \m{c\mapsto\alpha_S(c)(g)}.
\item For \m{x\in\WC\otimes Y_S} and \m{y\in\WC\otimes Y_{S'}}:
$n\cdot\pairW{\Pi}(x,y)=\pairW{\Pi}(n\cdot x,n\cdot y)$.
\end{enumerate}
\end{lemma}
\leanline{\lean{wittFstsmul}\sep \lean{smulwittPairing}}
\begin{proof}
(a) On \m{c_0\otimes(y,c)} both sides are $n c_0\otimes(gy+\alpha_S(c)(g))$. (b) The pairing modulo~$p^{M}$
satisfies $\langle gy,gy'\rangle=\chi(g)\langle y,y'\rangle$, and $\varphi$ acts on the first factors only.
\end{proof}

In particular, if $(1\otimes\projC{S})(x)=p^{e_1}(1\otimes c)$, then for \m{g\in\Gamma}
\begin{equation}\label{e:comp}
\projY{S}(gx)=g\,\projY{S}(x)+p^{e_1}\bigl(1\otimes\alpha_S(c)(g)\bigr),\qquad
\projY{S}(\varphi x)=\varphi\,\projY{S}(x).
\end{equation}

\begin{proof}[Proof of Theorem A]
Fix a topology on~$\Gamma$ for which $j_\Gamma$ is an open embedding. Write \m{a=\alpha_S(c)},
\m{a'=\alpha_{S'}(c')}, \m{\xi=\projY{S}(x)}, \m{\xi'=\projY{S'}(x')}, \m{\xi_w=\projY{S}(w)},
\m{\mathsf{a}_g=1\otimes a(g)}, \m{\mathsf{a}'_g=1\otimes a'(g)}, and
$\mathsf{B}(y,y')=p^{e_3}\pairW{\Pi}(y,y')$. The maps $a$, $a'$ are cocycles:
$\mathsf{a}_{gh}=\mathsf{a}_g+g\mathsf{a}_h$. By~\eqref{e:comp}, for \m{g,h\in\Gamma'},
\[\projY{S}(gx-x)=g\xi-\xi+p^{e_1}\mathsf{a}_g,\qquad
\projY{S'}\bigl(g(hx'-x')\bigr)=g(h\xi'-\xi')+p^{e_2}g\mathsf{a}'_h ,\]
and similarly for $w$ and for~$\varphi x'$, whose components in \m{\WC\otimes C_S}, \m{\WC\otimes C_{S'}} are again
\m{p^{e_1}(1\otimes c)} and \m{p^{e_2}(1\otimes c')}.

\emph{The cochain.} By Input~\ref{j:asw}(a) there is \m{y_0\in\mathbb{M}_M} with
$\varphi y_0-y_0=\mathsf{B}(\xi_w,\varphi\xi'-\xi')$. For \m{g\in\Gamma'} put
\[D_g=p^{e_1}\mathsf{B}(\mathsf{a}_g,g\xi')-\mathsf{B}(\xi,g\xi'-\xi')-p^{e_2}\mathsf{B}(\xi,\mathsf{a}'_g)
+(gy_0-y_0)\ \in\ \mathbb{M}_M .\]
A direct computation with Lemma~\ref{l:comp}(b) and~\eqref{e:comp} gives
\begin{align*}\varphi D_g-D_g={}&\mathsf{B}\bigl(\projY{S}(gw-w),\projY{S'}(g(\varphi x'-x'))\bigr)
-\mathsf{B}\bigl(\projY{S}(\varphi x-w),\projY{S'}(g\varphi x'-\varphi x')\bigr)\\
&-\mathsf{B}\bigl(\projY{S}(w-x),\projY{S'}(gx'-x')\bigr):\end{align*}
on the right the terms with~$\xi_w$ combine to
$\mathsf{B}(g\xi_w,g(\varphi\xi'-\xi'))-\mathsf{B}(\xi_w,\varphi\xi'-\xi')=(g-1)(\varphi y_0-y_0)$, the terms
with \m{\mathsf{a}_g} to $p^{e_1}\mathsf{B}(\mathsf{a}_g,g(\varphi\xi'-\xi'))$, those with
\m{\mathsf{a}'_g} to $-p^{e_2}\mathsf{B}(\varphi\xi-\xi,\mathsf{a}'_g)$, and the rest to
$-\mathsf{B}(\varphi\xi,g\varphi\xi'-\varphi\xi')+\mathsf{B}(\xi,g\xi'-\xi')$. The three terms on the right
lie in~$\mathfrak{s}_M$ by the last three hypotheses. By Input~\ref{j:asw}(b) there is
\m{k(g)\in\mathbb{Z}/p^{M}} with $D_g-1\otimes k(g)\in\mathfrak{s}_M$.

\emph{The coboundary.} For \m{g,h\in\Gamma'} the cocycle relations give
\[gD_h-D_{gh}+D_g=-p^{e_1}\mathsf{B}\bigl(\mathsf{a}_g,g(h\xi'-\xi')\bigr)
-\mathsf{B}\bigl(g\xi-\xi,g(h\xi'-\xi')\bigr)-p^{e_2}\mathsf{B}\bigl(g\xi-\xi,g\mathsf{a}'_h\bigr),\]
the terms with~$y_0$ cancelling. On the other hand the first hypothesis says that
\[\mathsf{B}\bigl(g\xi-\xi+p^{e_1}\mathsf{a}_g,\ g(h\xi'-\xi')+p^{e_2}g\mathsf{a}'_h\bigr)\in\mathfrak{s}_M .\]
Adding the two, and using that $\mathfrak{s}_M$ is stable under~$\Gamma$:
\[p^{e_1+e_2+e_3}\bigl(1\otimes\langle a(g),g\,a'(h)\rangle\bigr)-1\otimes\bigl(\chi(g)k(h)-k(gh)+k(g)\bigr)
\in\mathfrak{s}_M ,\]
so by Input~\ref{j:asw}(c), with \m{r=p^{e-e_1-e_2-e_3}},
\begin{equation}\label{e:cob}
p^{e}\langle a(g),g\,a'(h)\rangle=r\,\bigl(\chi(g)k(h)-k(gh)+k(g)\bigr)\qquad(g,h\in\Gamma').
\end{equation}

\emph{Continuity.} By Input~\ref{j:cont}, \m{a(g)\in p^{M}Y_S} and \m{a'(g)\in p^{M}Y_{S'}} for $g$ near~$1$,
and \m{gt-t\in p^{M}Y_{S'}} near~$1$ for each \m{t\in Y_{S'}}. Hence the function
$\mathsf{c}(g,h)=\mathsf{e}_M(p^{e}\langle a(g),g\,a'(h)\rangle)$ on \m{\Gamma\times\Gamma} is locally constant; it
is a $2$-cocycle, being a cup product of cocycles. For $h$ near~$1$ we have \m{D_h\in\mathfrak{s}_M}: the terms
with \m{\mathsf{a}_h}, \m{\mathsf{a}'_h} vanish because $p^{M}$ kills~$\mathbb{M}_M$, and
\m{\mathsf{B}(\xi,h\xi'-\xi')} and \m{hy_0-y_0} lie in~$\mathfrak{s}_M$ by Input~\ref{j:asw}(d), applied
to~$Y_{S'}$ and to \m{\mathbb{Z}/p^{M}(1)}. So \m{k(h)=0} for $h$ near~$1$, and then~\eqref{e:cob} shows
\m{k(gh)=k(g)}: the cochain~$k$ on~$\Gamma'$ is locally constant.

\emph{Conclusion.} The subgroup~$\Gamma'$ is open, because it contains the elements near~$1$. By~\eqref{e:cob}
the restriction of~$\mathsf{c}$ to~$\Gamma'$ is the coboundary of the continuous cochain
\m{\mathsf{e}_M\circ(rk)}. By Input~\ref{j:res}, the index of~$\Gamma'$ kills the class of~$\mathsf{c}$. This
class is also killed by~$p^{M}$, and the index is prime to~$p$; so the class is trivial, which is the assertion.
\end{proof}

\sect{Points: formal properties, strictness and deepening}{s:points}
Throughout, the setting is Setting~\ref{set:base}; $S$ is a side over~$\cO$.

\begin{lemma}[small facts]\label{l:facts}
Let $i$ be a letter and $o$, $o_1$, $o_2$, $o_3$ exponent functions.
\begin{enumerate}[label=\textup{(\alph*)},leftmargin=2.3em,itemsep=2pt]
\item For \m{y\in\WC\otimes Y_S}: $\projY{S}\bigl((1\otimes\inclY{S})(y)\bigr)=y$ and
$(1\otimes\projC{S})\bigl((1\otimes\inclY{S})(y)\bigr)=0$.
\item Let \m{e\ge0} be an integer, $\mathsf{A}$ an additive map with \m{\Reaches{S}{e}{\mathsf{A}}}, and $N$ a
$\Zp$-submodule of \m{\Mloc(\extL{S})} with $\locKP{\kerX{S}{i}{o}}\subseteq N$. If \m{x\in N} and
\m{x\in\deepE{S}{i}{o}}, then \m{\mathsf{A}(x)\in N}.
\item $\locKP{\kerX{S}{i}{o}}\subseteq\deepE{S}{i}{o}$.
\item Let \m{c\in C_S}, \m{e\ge0}, $x\in\ptsOver{S}{c}{i}{e}{o_1}{o}$ and \m{g\in\Gamma'}. Then
$(z\tau)^{-1}(gx-x)\in\kerX{S}{i}{o}$.
\item Let \m{c\in C_S}, \m{e\ge0}, $x\in\ptsOver{S}{c}{i}{e}{o_1}{o}$ and
$w\in\ptsOver{S}{c}{i}{e}{o_2}{o_3}$. Then $w-x\in\kerX{S}{i}{\min(o_1,o_2)}$.
\end{enumerate}
\end{lemma}
\leanline{\lean{wittFstinclKer}\sep \lean{memofstable}\sep \lean{kerWithExponentsledeepElements}\sep
\lean{smulsubmemkerWithExponents}\sep \lean{submemkerWithExponentsofpointsOver}\sep \lean{stepBackzero} (the
identity \m{i^{p^{0}}=i})}
\begin{proof}
(a) is clear. (b) $\mathsf{A}(x)=p^{e}x+(\mathsf{A}(x)-p^{e}x)$. (c) Such an element lies over
\m{0=\constAt{S}{i}{0}}, and its Galois differences keep the exponents (Input~\ref{q:exp}). (d) By
Input~\ref{q:ztau}; the image in \m{\WC\otimes C_S} is~$0$ because $g$ fixes \m{\constAt{S}{i}{c}}. (e) By
Input~\ref{q:exp}, and both points lie over the same element.
\end{proof}

\begin{lemma}[safety is checked letter by letter]\label{l:sum}
In Setting~\ref{set:pair}, let $S$, $S'$ be sides, $\Pi$ a pairing of $S$ and~$S'$, \m{\kappa\ge0} an integer,
and for every letter~$j$ let \m{u_j\in(\WC\otimes\extL{S})_j} and \m{v_j\in(\WC\otimes\extL{S'})_j}. If
\m{\Safe{\Pi}{\kappa}{u_j}{v_j}} for every letter~$j$, then
$\Safe{\Pi}{\kappa}{\sum_ju_j}{\sum_jv_j}$.
\end{lemma}
\leanline{\lean{isSafesum}}
\begin{proof}
The form~$\psi$ is $\cO$-bilinear, so the pairing modulo~$p^{M}$ satisfies
$\langle\omega y,y'\rangle=\langle y,\omega y'\rangle$. For letters \m{j\ne j'} this gives
$([j]-[j'])\,\pairW{\Pi}(\projY{S}(u_j),\projY{S'}(v_{j'}))=0$, and \m{[j]-[j']} is a unit
(Input~\ref{j:mod}). So the cross terms vanish.
\end{proof}

\begin{lemma}[extending a period to the extension]\label{l:extend}
Let $M$ be a $B^{+}_{\mathrm{dR}}$-submodule of~$\BdR$ and \m{\mu\in\Per(Y_S)} with \m{\mu(1\otimes t)\in M} for
all \m{t\in Y_S}. Then there is \m{\mu'\in\Per(\extL{S})} with $\mu'(\Ainf\otimes\extL{S})\subseteq M$ and
$\mu'\circ(1\otimes\inclY{S})=\mu$.
\end{lemma}
\leanline{\lean{existsextendform}}
\begin{proof}
Let $\Lambda\colon B_{\max}\otimes Y_S\to\BdR$ be the $B_{\max}$-linear map with
$\Lambda(1\otimes t)=\mu(1\otimes t)$; it commutes with~$\Gamma'$. For \m{c\in C_S} choose
\m{y_c\in B_{\max}\otimes Y_S} with $1\otimes\alpha_S(c)(g)=gy_c-y_c$ for all~$g$. Then
$g\Lambda(y_c)-\Lambda(y_c)=\mu(1\otimes\alpha_S(c)(g))\in M$ for \m{g\in\Gamma'}, so by Input~\ref{j:tate}
there is \m{a_c\in F} with \m{\Lambda(y_c)-a_c\in M}. Take a $\Zp$-basis \m{(e_k)} of~$C_S$, let
$\mu_C\colon\WC\otimes C_S\to\BdR$ be $\Ainf$-linear with $\mu_C(1\otimes e_k)=\Lambda(y_{e_k})-a_{e_k}$
(Input~\ref{j:types}), and put $\mu'=\mu\circ\projY{S}+\mu_C\circ(1\otimes\projC{S})$. Its values on
\m{\Ainf\otimes\extL{S}} lie in~$M$. It is a period: for \m{g\in\Gamma'},
$\mu'(1\otimes g(t,c))=g\mu(1\otimes t)+\mu(1\otimes\alpha_S(c)(g))+\mu_C(1\otimes c)$, and
$\mu(1\otimes\alpha_S(c)(g))=(g-1)\mu_C(1\otimes c)$ by linearity in~$c$, the~$a_{e_k}$ being fixed
by~$\Gamma'$.
\end{proof}

\begin{lemma}[strictness]\label{l:strict}
Let $i$ be a letter, $o$ an exponent function and \m{w\in\kerX{S}{i}{o}}. Then there is
\m{y\in\wExp{Y_S}{i}{o}} with \m{(1\otimes\inclY{S})(y)=w}; and if \m{w\in\Mloc(\extL{S})}, then
\m{y\in\Mloc(Y_S)}.
\end{lemma}
\leanline{\lean{existsinclKereqofprojTriveqzero}}
\begin{proof}
The sequence $0\to\WC\otimes Y_S\to\WC\otimes\extL{S}\to\WC\otimes C_S\to0$ is exact, so
\m{w=(1\otimes\inclY{S})(y)} with \m{y=\projY{S}(w)}, and $y$ lies in the part at~$i$. Some product~$\pi$ of
the~$\cE_\nu$ has \m{\pi w\in\Minf(\extL{S})}; then $\tau^{a}\pi y\in\Ainf\otimes Y_S$ for some~$a$, and for
\m{\mu\in\Per(Y_S)}, extended to~$\mu'$ by Lemma~\ref{l:extend} with \m{M=\BdR}, and every \m{b\ge0},
$\mu(\varphi^{b+1}\pi y)=\mu'(\varphi^{b+1}\pi w)\in\Fil^{0}\BdR$. So \m{y\in\Mlocinf(Y_S)}. For the
exponents let \m{j=i^{p^{-\nu}}}, \m{u\in U_{Y_S}(j)} and \m{\mu\in\Per_{j,u}(Y_S)}. The period
\m{\mu\circ\pr_j} takes values in \m{\Fil^{-u}\BdR} on \m{1\otimes Y_S}; its extension~$\mu'$ of
Lemma~\ref{l:extend} lies in \m{\Per_{j,u}(\extL{S})}, and
$\mu(\varphi^{-\nu}y)=\mu'(\varphi^{-\nu}w)\in\Fil^{o(\nu,u)}\BdR$. Finally, if \m{w\in\Mloc(\extL{S})}, then
$(1\otimes\inclY{S})(\ret_{Y_S}(y))=\ret_{\extL{S}}(w)=w$ by Input~\ref{q:ret}, so \m{\ret_{Y_S}(y)=y} and
\m{y\in\Mloc(Y_S)}.
\end{proof}

\begin{lemma}[a point from a cokernel killer]\label{l:first}
Let $i$ be a letter, \m{\kappa,\rho\ge0} integers with \m{\Kills{S}{i}{\kappa}{\rho}}, and \m{c\in C_S}. Then
$\ptsOver{S}{c}{i}{\kappa}{\raisedE{-\rho}{\cdot}{\cdot}}{\raisedE{-\rho}{\cdot}{\cdot}}$ is not empty.
\end{lemma}
\leanline{\lean{pointsOvernonemptyofkillsCokernel}}
\begin{proof}
By Input~\ref{q:kill} there is \m{x_1\in\locLP{0}{i}(\extL{S})} over
$p^{\kappa}\cE^{\rho}\constAt{S}{i}{c}$. Then \m{x=\cE^{-\rho}x_1} lies in \m{\Mloc(\extL{S})}, over
\m{p^{\kappa}\constAt{S}{i}{c}}, with the exponents \m{\std_0} lowered by~$\rho$ at \m{\nu=0}
(Inputs \ref{q:loc} and~\ref{q:exp}); its Galois differences have the same exponents.
\end{proof}

\begin{proposition}[deepening along a set of steps]\label{p:reach}
Let $I$ be a finite set of pairs \m{(\nu,n)} of an integer \m{\nu\ge0} and an integer~$n$. Then there is an
additive map $\mathsf{A}\colon\Mloc(\extL{S})\to\Mloc(\extL{S})$ with
$\Reaches{S}{\sum_{(\nu,n)\in I}\cost{\nu}{n}}{\mathsf{A}}$ such that for every letter~$i$, all exponent
functions $o$, $o'$ and every \m{x\in\deepE{S}{i}{o}}: if $I$ contains every \m{(\nu,n)} with
\m{o(\nu,u)\le n<o'(\nu,u)} for some \m{u\in U_{Y_S}(i^{p^{-\nu}})}, then $\mathsf{A}(x)\in\deepE{S}{i}{o'}$.
\end{proposition}
\leanline{\lean{existsreachesdeep}}
\begin{proof}
By induction on~$I$, adding at each stage a pair whose second component is maximal. For \m{I=\emptyset} take the
identity: then \m{o'\le o} on the jump sets. Let \m{I=I_0\cup\{(\nu_0,n_0)\}} with \m{n\le n_0} for all
\m{(\nu,n)\in I_0}, and let $\mathsf{A}_0$ serve for~$I_0$. Take \m{t_1,\dots,t_m} and~$y$ from
Input~\ref{q:deep} for \m{(\nu_0,n_0)} and the lattice datum \m{\extL{S}}, and put
\[\mathsf{B}(x)=p^{\cost{\nu_0}{n_0}}x+\sum_k\ret_{\extL{S}}\bigl(y\,(z\tau)^{-1}(t_kx-x)\bigr),\qquad
\mathsf{A}=\mathsf{B}\circ\mathsf{A}_0 .\]
\emph{Cost.} For \m{x\in\deepE{S}{i}{o}} each \m{(z\tau)^{-1}(t_kx-x)} lies in \m{\wExp{\extL{S}}{i}{o}}
(Input~\ref{q:ztau}) and maps to~$0$ in \m{\WC\otimes C_S}, and $\ret$ keeps both properties
(Input~\ref{q:ret}); so $\mathsf{B}$ has cost \m{\cost{\nu_0}{n_0}}. Costs add under composition, by
Lemma~\ref{l:facts}(b) and~(c).

\emph{Depth.} Let \m{x\in\deepE{S}{i}{o}} and let $I$ contain the steps from $o$ to~$o'$. Let $o''$ be $o'$
lowered to~$n_0$ at the places \m{(\nu_0,u)} where \m{o'(\nu_0,u)=n_0+1}. By the maximality of~$n_0$, $I_0$
contains the steps from $o$ to~$o''$, so $x_1=\mathsf{A}_0(x)\in\deepE{S}{i}{o''}$, and
$\mathsf{B}(x_1)\in\deepE{S}{i}{o''}$ by Lemma~\ref{l:facts}(b). By Input~\ref{q:exp} it remains to check the
Galois differences of \m{\mathsf{B}(x_1)} at the places where \m{o''<o'}: so let \m{u} with
\m{o'(\nu_0,u)=n_0+1} and $\mu\in\Per_{i^{p^{-\nu_0}},u}(\extL{S})$. The value
\m{V=\mu(\varphi^{-\nu_0}x_1)} has \m{gV-V\in\Fil^{n_0}\BdR} for all \m{g\in\Gamma'}, so
\m{V-c\in\Fil^{n_0}\BdR} for some \m{c\in F} (Input~\ref{j:tate}). By Input~\ref{q:deep}, the value of~$\mu$ at
\m{\varphi^{-\nu_0}} of $p^{\cost{\nu_0}{n_0}}x_1+\sum_ky(z\tau)^{-1}(t_kx_1-x_1)$ is congruent to some
\m{c'\in F} modulo \m{\Fil^{n_0+1}\BdR}, and by Input~\ref{q:ret} the same holds for its
retraction~\m{\mathsf{B}(x_1)}. As $c'$ is fixed by~$\Gamma'$, the Galois differences of that value lie in
\m{\Fil^{n_0+1}\BdR}.
\end{proof}

\begin{lemma}[points after deepening; Frobenius]\label{l:move}
Let \m{c\in C_S}, $i$ a letter, \m{e,e',l\ge0} integers and $o$, $o'$, $o_1$ exponent functions.
\begin{enumerate}[label=\textup{(\alph*)},leftmargin=2.3em,itemsep=2pt]
\item Let $x\in\ptsOver{S}{c}{i}{e}{o_1}{o}$ and let $I$ be a finite set of pairs \m{(\nu,n)} that
contains every \m{(\nu,n)} with \m{o(\nu,u)\le n<o'(\nu,u)} for some \m{u\in U_{Y_S}(i^{p^{-\nu}})}. If
$e+\sum_{(\nu,n)\in I}\cost{\nu}{n}\le e'$, then
$\ptsOver{S}{c}{i}{e'}{\impE{o_1}{o'}{\cdot}{\cdot}}{o'}$ is not empty.
\item If $w\in\ptsOver{S}{c}{i^{p^{-l}}}{e}{o}{o'}$, then
$\varphi^{l}w\in\ptsOver{S}{c}{i}{e}{\frobX{l}{0}{o}{\cdot}{\cdot}}{\frobX{l}{1}{o'}{\cdot}{\cdot}}$.
\item Let \m{n\ge0} be an integer and $\beta$ an integer with \m{\NoMult{p}{n}{\beta}}, and let
$o^{\mathrm{lo}}$, $o^{\mathrm{hi}}$ be exponent functions. Suppose that
$\ptsOver{S}{c}{i}{\kappa}{o^{\mathrm{lo}}}{o}$ is not empty, for an integer \m{\kappa\ge0}, and that
for all \m{\nu\ge0}, \m{u\in U_{Y_S}(i^{p^{-\nu}})} and integers~$a$ with
$o(\nu,u)\le a<o^{\mathrm{hi}}(\nu,u)$ one has \m{\nu=0} and
$a\in\twoWin{n}{\beta}\cup\twoWin{n}{-\beta}$. Then
$\ptsOver{S}{c}{i}{\kappa+(8n+4)\fd}{\impE{o^{\mathrm{lo}}}{o^{\mathrm{hi}}}{\cdot}{\cdot}}{o^{\mathrm{hi}}}$ is
not empty.
\end{enumerate}
\end{lemma}
\leanline{\lean{pointsOvernonemptyofcosts}\sep \lean{phipowsmulmempointsOver}\sep
\lean{pointsOvernonemptyofnoMultipleNear}}
\begin{proof}
(a) Let $\mathsf{A}$ be the map of Proposition~\ref{p:reach} for~$I$, of cost~$\epsilon$, and
\m{x_2=p^{e'-e-\epsilon}x}. It lies over a constant, in \m{\deepE{S}{i}{o}} and in \m{\deepE{S}{i}{o_1}}. So
\m{\mathsf{A}(x_2)} lies over \m{p^{e'}\constAt{S}{i}{c}}, in \m{\wExp{\extL{S}}{i}{o_1}}
(Lemma~\ref{l:facts}(b)) and in \m{\deepE{S}{i}{o'}}; Input~\ref{q:deep} improves its exponents.

(b) \m{\Mloc} is stable under~$\varphi$, $\varphi^{l}\constAt{S}{i^{p^{-l}}}{c}=\constAt{S}{i}{c}$, and the
exponents shift by Input~\ref{q:B}. For the depth, $(z\tau)^{-1}(gw-w)\in\wExp{\extL{S}}{i^{p^{-l}}}{o'}$
by Lemma~\ref{l:facts}(d), so its image under~$\varphi^{l}$ has the exponents
\m{\frobX{l}{0}{o'}{\cdot}{\cdot}}; and
$g\varphi^{l}w-\varphi^{l}w=\varphi^{l}(z\tau)\cdot\varphi^{l}\bigl((z\tau)^{-1}(gw-w)\bigr)$, where
\m{\varphi^{l}(\tau)} is divisible by~$\cE_\nu$ for every \m{\nu<l}, which raises the exponents at these~$\nu$
by~$1$.

(c) Apply (a) with $I=\{(0,a): a\in\twoWin{n}{\beta}\cup\twoWin{n}{-\beta}\}$. This set has at most
\m{8n+4} elements, and \m{\cost{0}{a}=v_p(a)+\fd=\fd} for each of them, because $p$ divides none of them
except possibly~$0$.
\end{proof}

\sect{Safety from exponents}{s:exponents}
Throughout, the setting is Setting~\ref{set:pair}; $S$, $S'$ are sides over~$\cO$ and $\Pi$ a pairing of $S$
and~$S'$ modulo~$p^{M}$. Write $\mathsf{m}\colon\cO\times Y_S\to Y_S$ for the scalar multiplication,
$\mathsf{t}\colon Y_S\times Y_{S'}\to Y_S\otimes_{\cO}Y_{S'}$ for the tensor map, and
$\bar{\operatorname{tr}}\colon\cO(1)\to\mathbb{Z}/p^{M}(1)$ for $\operatorname{tr}$ followed by reduction.

\begin{lemma}[the pairing through the third side]\label{l:third}
Let $i$ be a letter. Then:
\begin{enumerate}[label=\textup{(\alph*)},leftmargin=2.3em,itemsep=2pt]
\item $\wPair{\mathsf{m}}{\constAt{\third{\Pi}}{i}{1}}{w}=\pr_i(w)$ for every \m{w\in\WC\otimes Y_S};
\item for every $n\in\WC\otimes\extL{\third{\Pi}}$, \m{w\in\WC\otimes Y_S} and \m{y\in\WC\otimes Y_{S'}},
\[(1\otimes\bar{\operatorname{tr}})\Bigl(\wPair{\mathsf{p}_\Pi}{n}{\wPair{\mathsf{t}}{w}{y}}\Bigr)
=\pairW{\Pi}\Bigl(\wPair{\mathsf{m}}{(1\otimes\projC{\third{\Pi}})(n)}{w},\ y\Bigr).\]
\end{enumerate}
\end{lemma}
\leanline{\lean{wittPairingconstantAtandthirdPairing}\sep the formal statement binds the extension of the scalar
multiplication by a local definition before the two assertions}
\begin{proof}
Both are identities of $\WC$-multilinear maps; check them on elementary tensors. (a) is
Input~\ref{q:C}, last sentence. For~(b), with \m{n=1\otimes(\lambda,a)}:
$\bar{\operatorname{tr}}(p^{M}\lambda(t\otimes t')+a\psi(t\otimes t'))=\langle at,t'\rangle$.
\end{proof}

\begin{proposition}[safety from exponents]\label{p:safe}
Let $i$ be a letter, \m{l,\kappa\ge0} integers and $o$, $o'$, $o_2$, $o_3$ exponent functions. Let
$n\in\ptsOver{\third{\Pi}}{1}{i^{p^{-l}}}{\kappa}{o_2}{o_3}$, and suppose that \m{u\mapsto o_2(\nu,u)} is
non-increasing for every \m{\nu\ge0}. Let \m{w\in\kerX{S}{i}{o}} and \m{y\in\kerX{S'}{i}{o'}}, and suppose that
\begin{equation}\label{e:pos}
0<\frobX{l}{0}{o_2}{\nu}{-u-u'-1}+o(\nu,u)+o'(\nu,u')
\end{equation}
for all \m{\nu\ge0}, \m{u\in U_{Y_S}(i^{p^{-\nu}})} and \m{u'\in U_{Y_{S'}}(i^{p^{-\nu}})}. Then
$\Safe{\Pi}{\kappa}{z\tau\cdot w}{y}$ and $\Safe{\Pi}{\kappa}{w}{z\tau\cdot y}$.
\end{proposition}
\leanline{\lean{isSafeofexponents}}
\begin{proof}
By Lemma~\ref{l:strict}, \m{w=(1\otimes\inclY{S})(w_0)} and \m{y=(1\otimes\inclY{S'})(y_0)} with
\m{w_0\in\wExp{Y_S}{i}{o}}, \m{y_0\in\wExp{Y_{S'}}{i}{o'}}. By Lemma~\ref{l:move}(b),
\m{n_1=\varphi^{l}n} is a point of~\m{\third{\Pi}} at~$i$ over \m{p^{\kappa}\cdot1}, with the exponents
$o_4=\frobX{l}{0}{o_2}{\cdot}{\cdot}$, again non-increasing in~$u$.

\emph{Exponents.} Apply Input~\ref{q:C} to the tensor map~$\mathsf{t}$:
\[\wPair{\mathsf{t}}{w_0}{y_0}\in\wExp{\tensL{S}{S'}}{i}{o_5}\qquad\text{with}\qquad o_5(\nu,v)=1-o_4(\nu,-v-1).\]
The condition is $o_5(\nu,v)\le o(\nu,u)+o'(\nu,u')$ for \m{u+u'\le v}; it follows from~\eqref{e:pos},
because $o_4$ is non-increasing in its second argument and \m{-v-1\le-u-u'-1}. Apply it again to the third pairing~$\mathsf{p}_\Pi$, with values
in~\m{\tateL{\cO}}:
\[q=\wPair{\mathsf{p}_\Pi}{n_1}{\wPair{\mathsf{t}}{w_0}{y_0}}\ \in\ \wExp{\tateL{\cO}}{i}{\std_0};\]
the condition is $\std_0(\nu,-1)=1\le o_5(\nu,v)+o_4(\nu,u_3)$ for jumps~$u_3$ of the third side with
\m{u_3+v\le-1}, which holds because \m{o_4(\nu,u_3)\ge o_4(\nu,-v-1)}. By Input~\ref{q:kill} with the first
power of~$\chi$, \m{\tau q\in\Ainf\otimes\cO(1)}.

\emph{The pairing.} By Lemma~\ref{l:third} and $(1\otimes\projC{\third{\Pi}})(n_1)=p^{\kappa}\constAt{\third{\Pi}}{i}{1}$,
\[(1\otimes\bar{\operatorname{tr}})(q)=p^{\kappa}\,\pairW{\Pi}(\pr_i(w_0),y_0)=p^{\kappa}\,\pairW{\Pi}(w_0,y_0).\]
Hence $p^{\kappa}\pairW{\Pi}(z\tau w_0,y_0)=z\cdot(1\otimes\bar{\operatorname{tr}})(\tau q)$. The element
$(1\otimes\bar{\operatorname{tr}})(\tau q)$ is an $\Ainf$-combination of elements \m{1\otimes k}, and
\m{z\in W(\mathfrak{m}^\flat)} (Input~\ref{j:mod}); so the product lies in~$\mathfrak{s}_M$. As
\m{\projY{S}(w)=w_0} and \m{\projY{S'}(y)=y_0} (Lemma~\ref{l:facts}(a)) and the pairing is bilinear, both
assertions follow.
\end{proof}

\begin{lemma}[safety of a Galois difference against a translate]\label{l:safediff}
In the situation of Proposition~\ref{p:safe}, with the point~$n$, the functions \m{o,o',o_2,o_3},
\m{y\in\kerX{S'}{i}{o'}} and~\eqref{e:pos} as there, let \m{c\in C_S}, \m{e\ge0} an integer, $o_1$ an
exponent function and $x\in\ptsOver{S}{c}{i}{e}{o_1}{o}$. Then
$\Safe{\Pi}{\kappa}{gx-x}{g\cdot y}$ for every \m{g\in\Gamma'}.
\end{lemma}
\leanline{\lean{isSafesmulsub}}
\begin{proof}
$(z\tau)^{-1}(gx-x)\in\kerX{S}{i}{o}$ by Lemma~\ref{l:facts}(d), and \m{gy=(gy-y)+y\in\kerX{S'}{i}{o'}} by
Input~\ref{q:exp}. Apply Proposition~\ref{p:safe} to these two elements.
\end{proof}

\sect{The walk}{s:walk}
\begin{proof}[Proof of Theorem B]
For each letter~$j$ with the property~$\mathsf{cut}$ (a \emph{cut}) choose: from~(i), an exponent function~$d_j$,
an integer~$r_j$ and a point $x_j\in\ptsOver{S}{c}{j}{e}{o(j)}{d_j}$; from~(ii), $d'_j$, $r'_j$ and
$x'_j\in\ptsOver{S'}{c'}{j}{e'}{o'(j)}{d'_j}$; from~(iii), a point
$n_j\in\ptsOver{\third{\Pi}}{1}{j}{\kappa}{z(j)}{z'(j)}$. For a letter~$j$ let \m{\mathsf{l}=\mathsf{l}(j)}
and \m{j^-=j^{p^{-\mathsf{l}}}}, the nearest cut strictly before~$j$ on its Frobenius orbit, and put
\[X^{\bot}_j=\varphi^{\mathsf{l}}x_{j^-},\qquad X^{\top}_j=x_j\ \text{ if $j$ is a cut},\quad
X^{\top}_j=X^{\bot}_j\ \text{ otherwise},\]
\[X^{\mathsf{Q}}_j=x_j\ \text{ if $j$ is a cut and }\mathsf{Q}(j),\quad X^{\mathsf{Q}}_j=X^{\bot}_j\
\text{ otherwise},\]
and define \m{X'^{\bot}_j}, \m{X'^{\top}_j} in the same way from the~$x'_j$. By Lemma~\ref{l:move}(b) these
are points at~$j$ over \m{p^{e}c}, respectively \m{p^{e'}c'}; a point moved from~$j^-$ has the exponents
$\frobX{\mathsf{l}}{0}{o(j^-)}{\cdot}{\cdot}$ and the depth $\frobX{\mathsf{l}}{1}{d_{j^-}}{\cdot}{\cdot}$.
Let \m{X^{\top}=\sum_jX^{\top}_j}, and likewise \m{X^{\bot}}, \m{X^{\mathsf{Q}}}, \m{X'^{\top}},
\m{X'^{\bot}}. The component of \m{\varphi X^{\top}} at~$j$ is \m{\varphi(X^{\top}_{j^{1/p}})}, and the nearest cut at or
before~\m{j^{p^{-1}}} is the nearest cut strictly before~$j$; hence \m{\varphi X^{\top}=X^{\bot}} and
\m{\varphi X'^{\top}=X'^{\bot}}.

We apply Theorem~A to \m{x=X^{\top}}, \m{w=X^{\mathsf{Q}}}, \m{x'=X'^{\top}}, with \m{(e_1,e_2,e_3)=(e,e',\kappa)}.
The three elements lie over \m{p^{e}(1\otimes c)}, \m{p^{e}(1\otimes c)}, \m{p^{e'}(1\otimes c')}, because
$\sum_j\constAt{S}{j}{c}=1\otimes c$. By Lemma~\ref{l:sum} the four safety conditions may be checked on the
components at each letter~$j$, which lie in the parts at~$j$. Fix $j$ and \m{g\in\Gamma'}.

\emph{First condition:} $(gX^{\top}_j-X^{\top}_j,\ g(hX'^{\top}_j-X'^{\top}_j))$. Let $j^{\circ}$ be the cut
from which \m{X^{\top}_j} was moved, at distance \m{a\ge0}. Apply Lemma~\ref{l:safediff} at~$j$ with the third
point~$n_{j^\circ}$ (and \m{l=a}), the point \m{X^{\top}_j} of depth $\frobX{a}{1}{d_{j^\circ}}{\cdot}{\cdot}$,
and the element \m{hX'^{\top}_j-X'^{\top}_j} of $\kerX{S'}{j}{\frobX{a}{0}{o'(j^\circ)}{\cdot}{\cdot}}$.
Condition~\eqref{e:pos} reads, for \m{\nu<a}, \m{0<0+1+0}; and for \m{\nu\ge a} it is clause~(a) of
Definition~\ref{d:deep} for~$S$ at~$j^\circ$, at the index \m{\nu-a}.

\emph{The other conditions} involve \m{X^{\mathsf{Q}}_j-X^{\top}_j}, \m{X'^{\bot}_j-X'^{\top}_j} and
\m{X^{\bot}_j-X^{\mathsf{Q}}_j}. If $j$ is not a cut, all three vanish and there is nothing to prove. Let $j$
be a cut.

\emph{Case \m{\mathsf{Q}(j)}.} Then $X^{\mathsf{Q}}_j=X^{\top}_j=x_j$; the second condition is trivial. The
third is $(gx_j-x_j,\ g(X'^{\bot}_j-x'_j))$: by Lemma~\ref{l:facts}(e),
$X'^{\bot}_j-x'_j\in\kerX{S'}{j}{\min(\frobX{\mathsf{l}}{0}{o'(j^-)}{\cdot}{\cdot},o'(j))}$, and
Lemma~\ref{l:safediff} applies with the point~$x_j$ of depth~$d_j$ and the third point at the cut
\m{j^{p^{-r_j}}}; condition~\eqref{e:pos} is the second inequality of clause~(b) for~$S$ at~$j$. The fourth
is $(X^{\bot}_j-x_j,\ gX'^{\bot}_j-X'^{\bot}_j)$: here
$X^{\bot}_j-x_j\in\kerX{S}{j}{\min(\frobX{\mathsf{l}}{0}{o(j^-)}{\cdot}{\cdot},o(j))}$ and
$(z\tau)^{-1}(gX'^{\bot}_j-X'^{\bot}_j)$ has the exponents
$\frobX{\mathsf{l}}{1}{d'_{j^-}}{\cdot}{\cdot}$; Proposition~\ref{p:safe} applies with the third point
at~$j^-$. Condition~\eqref{e:pos} holds for \m{\nu<\mathsf{l}} because $o(j)(\nu,u)\ge\std_{\mathsf{l}}(\nu,u)=0$
by the first inequality of clause~(b) for~$S$ at~$j$; and for \m{\nu\ge\mathsf{l}} by clause~(a) for~$S'$
at~$j^-$, which has the term $\min(-u,o(j^-)(\nu-\mathsf{l},u))$ on its right, together with
$o(j)(\nu,u)\ge\std_{\mathsf{l}}(\nu,u)=-u$.

\emph{Case not \m{\mathsf{Q}(j)}.} Then \m{X^{\mathsf{Q}}_j=X^{\bot}_j}; the fourth condition is trivial, and
the other two are the mirror images of the previous case with the roles of $S$ and~$S'$ exchanged: the second
is \m{(X^{\bot}_j-x_j,\ gx'_j-x'_j)}, treated by Proposition~\ref{p:safe} with the depth~$d'_j$ of~$x'_j$ and
the third point at the cut \m{j^{p^{-r'_j}}}, by the second inequality of clause~(b) for~$S'$ at~$j$; the third
is $(gX^{\bot}_j-X^{\bot}_j,\ g(X'^{\bot}_j-x'_j))$, treated by Lemma~\ref{l:safediff} with the third point
at~$j^-$, by clause~(a) for~$S$ at~$j^-$ and the first inequality of clause~(b) for~$S'$ at~$j$.
\end{proof}

\begin{proof}[Proof of Theorem C]
We apply Theorem~B with every letter a cut, so that \m{\mathsf{l}(i)=1} for all~$i$, with
$e=\kappa+(48(\rho+\eta+W)+4)\fd$ and an integer \m{e'\le e}, and with
$z(i)=z'(i)=\raisedE{-\rho}{\cdot}{\cdot}$, which is non-increasing in~$u$; the third points exist by
Lemma~\ref{l:first} and hypothesis~(i). Then \m{e+e'+\kappa\le H} by~(v). Put \m{n_0=6(\rho+\eta+W)} and
$D=\raisedE{2\rho+\eta}{\cdot}{\cdot}$. Write $i^-$ for~\m{i^{p^{-1}}}.

\emph{First alternative of~(iv).} Take \m{\mathsf{Q}(i)} to be \m{\varsigma(i)=1}. For $S$ at~$i$,
Lemma~\ref{l:first} gives a point over \m{p^{\kappa}c} with exponents and depth \m{\raisedE{-\rho}{\cdot}{\cdot}},
and Lemma~\ref{l:move}(c) deepens it to~$D$: the steps are at \m{\nu=0}, from \m{-\rho-u} to \m{2\rho+\eta-u},
and they lie in \m{[-n_0,n_0]} for \m{|u|\le\eta} and within~$n_0$ of \m{\pm\beta(i)} for~$u$ in the far window.
The result is a point over~$p^{e}c$ with depth~$D$ and exponents $o(i)=\impE{\raisedE{-\rho}{\cdot}{\cdot}}{D}{\cdot}{\cdot}$,
and \m{r=0}; the same for~$S'$, giving~$o'(i)$ and \m{e'=e}. Explicitly \m{o(i)(\nu,u)=-u} for \m{\nu\ge1}, and
\m{o(i)(0,u)} is \m{2\rho+\eta-u} for \m{u<-\rho} and \m{\min(2\rho+\eta-u,0)} for \m{u\ge-\rho}. It remains to check clauses (a) and~(b) of
Definition~\ref{d:deep}. At \m{\nu=0} the left side of~(a) is \m{-\rho-\eta-v-1}, which is smaller than~$-v$,
than \m{2\rho+\eta-v}, and than~$0$ when \m{v\ge-\rho}; at \m{\nu\ge1} it is \m{-v-1<-v=o'(i)(\nu,v)}. So (a)
holds for both sides. In~(b), for the side that leads at~$i$, the inequality \m{\std_1(0,u)=0\le o(i)(0,u)} asks
\m{u\le2\rho+\eta} when \m{u\ge-\rho}: the leading side has its far window on the negative side, so its jumps
are at most~$\eta$. The second inequality of~(b) asks in addition, at \m{\nu=0}, that \m{-\rho-\eta-v-1<0},
which holds because the jumps~$v$ of the other side are at least~$-\eta$; and at \m{\nu=1} that
\m{-v-1<o'(i^{p^{-1}})(0,v)}, which fails exactly when \m{-\rho\le v<0}: this is excluded by the
alternative. For \m{\nu\ge2} it reads \m{-v-1<-v}.

\emph{Second alternative of~(iv).} Now \m{\fd=0}. Take \m{\mathsf{Q}(i)} false for all~$i$, so that $S'$ leads
everywhere. The points of~$S$ and of the third side are as before. For~$S'$, whose jumps are all \m{\le0},
Input~\ref{q:kill} gives \m{\Kills{S'}{j}{0}{0}} for every letter~$j$, and Lemma~\ref{l:first} a point over~$c'$
at~$i^-$ with exponents and depth \m{\raisedE{0}{\cdot}{\cdot}}. Deepen it to~$D$ by Lemma~\ref{l:move}(c): the
steps are at \m{\nu=0} and lie in \m{[0,2\rho+2\eta)}. Apply~$\varphi$ (Lemma~\ref{l:move}(b)); the depth becomes
$1$ at \m{\nu=0} and \m{D(\nu-1,u)} at \m{\nu\ge1}. Deepen once more at \m{\nu=0}, from~$1$ to
\m{2\rho+\eta-u}. The result is a point over~$c'$ at~$i$, so that \m{e'=0}, with depth~$d'$ and
exponents~$o'(i)$ given by
\[\begin{gathered}d'(0,u)=d'(1,u)=2\rho+\eta-u,\qquad o'(i)(0,u)=0,\\
o'(i)(1,u)=2\rho+\eta-u\ \text{ for }u<0,\qquad o'(i)(1,0)=0,\end{gathered}\]
and \m{d'(\nu,u)=o'(i)(\nu,u)=-u} for \m{\nu\ge2}. Take \m{r'=1} at the letters with \m{\varsigma(i)=1} and
\m{r'=0} at the others.

\emph{Clause~(a) for~$S$.} The left side is \m{-\rho-\eta-v-1} at \m{\nu=0} and \m{-v-1} at \m{\nu\ge1}, as
before. At \m{\nu=0} the right side is \m{\min(-v,0)}, and \m{-\rho-\eta-v-1<0} because \m{v\ge-\eta}. At
\m{\nu=1} it is \m{\min(-v,o'(i)(1,v))}, which is larger than \m{-v-1} for \m{v<0} and equals~$0$ for
\m{v=0}. At \m{\nu\ge2} it is~$-v$.

\emph{Clause~(a) for~$S'$.} For \m{u'\in U_{Y_{S'}}}, \m{v\in U_{Y_S}} the left side is \m{-\rho-\eta-v-1} at
\m{\nu=0}, \m{-2\rho-\eta-v-1} at \m{\nu=1} and \m{-v-1} at \m{\nu\ge2}. It is smaller than~$-v$, and smaller
than \m{o(i)(\nu,v)}: at \m{\nu\ge1} this is~$-v$, and at \m{\nu=0} it was checked in the first alternative.

\emph{Clause~(b) for~$S'$.} First, \m{\std_1(\nu,u)\le o'(i)(\nu,u)}: at \m{\nu=0} both are~$0$, and at
\m{\nu\ge1} the left side is~$-u$. Second, the right side is the minimum of \m{o(i)(\nu,v)} and of the
function equal to~$0$ at \m{\nu=0} and to \m{o(i^-)(\nu-1,v)} at \m{\nu\ge1}. If \m{\varsigma(i)=-1}, the left
side is that of clause~(a). At \m{\nu=0} we need \m{-\rho-\eta-v-1<0}: the far window of~$S$ at~$i$ is on the
positive side, so \m{v\ge-\eta}. At \m{\nu=1} we need \m{-2\rho-\eta-v-1<o(i^-)(0,v)}: for \m{v<-\rho} the right
side is \m{2\rho+\eta-v}, and for \m{v\ge-\rho} it is \m{\min(2\rho+\eta-v,0)}, while the left side is
negative. At \m{\nu\ge2} we need \m{-v-1<-v}. If \m{\varsigma(i)=1}, the third point is moved from~$i^-$, and
the left side is \m{-(2\rho+\eta-u')} at \m{\nu=0}, \m{-\rho-\eta-v-1} at \m{\nu=1} and \m{-v-1} at
\m{\nu\ge2}. At \m{\nu=0} it is negative, because \m{\eta\ge1} and \m{u'\le0}, while
\m{\min(o(i)(0,v),0)=0}: the far window of~$S$ at~$i$ is on the negative side, so \m{v\le\eta} and
\m{o(i)(0,v)\ge0}. At \m{\nu=1} we need \m{-\rho-\eta-v-1<o(i^-)(0,v)}, which holds as in the previous case,
and at \m{\nu\ge2} again \m{-v-1<-v}.
\end{proof}

\sect{The long gap}{s:gap}
Throughout, the setting is Setting~\ref{set:base}.

\begin{proposition}[a relation among elements of the same weight]\label{p:auer}
Suppose that the Kummer chart is unramified. Let $L$ be a lattice datum, $o$ an exponent function, $i$ a letter,
\m{u\in\mathbb{Z}} and \m{n\ge0} an integer with \m{m_L(i,u)\le n}, where \m{m_L(i,u)} is the total
multiplicity of the jump~$u$ of~$L$ above~$i$. Let \m{\nu_0\ge0} be an integer and $j$ a letter with
\m{j^{p^{-\nu_0}}=i}, and suppose that \m{o(\nu_0,u)<o(\nu_0,v)} for every \m{v\in U_L(i)} with \m{v<u}. Let
\m{x_0,\dots,x_{hn}\in\Mloc(L)} with \m{x_k\in\wExp{L}{j}{o}} for all~$k$. Then there are \m{c_0,\dots,c_{hn}\in\Zp},
one of them equal to~$1$, with
\[\sum_kc_kx_k\ \in\ \locKP{\wExp{L}{j}{o_+}}+\locKP{\wExp{L}{j}{\raiseX{o}{\nu_0}{u}{\cdot}{\cdot}}},\]
where \m{o_+(\nu,v)=o(\nu,v)+1} for \m{\nu=0} and \m{o_+(\nu,v)=o(\nu,v)} otherwise.
\end{proposition}
\leanline{\lean{existscombofjumpMultiplicityle}}
\begin{proof}
Let \m{\mu_1,\dots,\mu_r}, \m{r\le n}, be the periods of Input~\ref{j:types} for \m{(i,u)}, and
\m{a=o(\nu_0,u)}. By Input~\ref{q:any} and the hypothesis on~$o$, $\mu_m(\varphi^{-\nu_0}x_k)\in\Fil^{a}\BdR$;
write it as \m{\pi^{a}t_{km}} with a fixed uniformizer~$\pi$ and \m{t_{km}\in B^{+}_{\mathrm{dR}}}, with
residue \m{\bar t_{km}\in C}. By Auerbach's lemma (Input~\ref{j:auer}) applied to the column space of the matrix
\m{(\bar t_{km})}, there are indices \m{s_1,\dots,s_q}, \m{q\le r}, and $\bar a_{kb}=\mathsf{z}_b(k)\in\mathcal{O}_C$
with $\bar t_{km}=\sum_b\bar a_{kb}\bar t_{s_bm}$ for all $k$,~$m$. Lift \m{\bar a_{kb}} to \m{a_{kb}\in\Ainf}
along~$\theta$ and put $y_k=x_k-\sum_b\varphi^{\nu_0}(a_{kb})x_{s_b}\in\wExp{L}{j}{o}$. Then
$\mu_m(\varphi^{-\nu_0}y_k)\in\Fil^{a+1}\BdR$ for all~$m$; since every period of type~$u$ at~$i$ is a
combination of the~$\mu_m$, with coefficients in~$\bar{\mathbb{Q}}_p$, modulo periods of type \m{u-1}, on which the value lies in \m{\Fil^{a+1}\BdR} by
Input~\ref{q:any}, we get $y_k\in\wExp{L}{j}{\raiseX{o}{\nu_0}{u}{\cdot}{\cdot}}$, and the same for
$\ret_L(y_k)=x_k-\sum_b\operatorname{co}_0(\varphi^{\nu_0}a_{kb})x_{s_b}$.
Write $\operatorname{co}_0(\varphi^{\nu_0}a_{kb})=w_{kb}+\cE y_{kb}$ with \m{w_{kb}\in\WW}
(Input~\ref{j:mod}). The \m{hn+1} vectors \m{(w_{kb})_b\in\WW^{q}} lie in a free $\Zp$-module of rank
\m{hq\le hn}, so they satisfy a non-trivial relation over~$\Zp$; dividing by a coefficient of least valuation
we get \m{c_k\in\Zp}, one of them~$1$, with \m{\sum_kc_kw_{kb}=0} for all~$b$. Then
\[\sum_kc_kx_k=\sum_kc_k\,\ret_L(y_k)+\cE\cdot\sum_{k,b}c_k\,y_{kb}\,x_{s_b},\]
and multiplication by~$\cE$ raises the exponents at \m{\nu=0} by~$1$.
\end{proof}

\begin{lemma}[coefficients of an expansion in powers of~$\cE$]\label{l:coeff}
Let $L$ be a lattice datum, \m{l>0} and \m{A\ge0} integers, $i$ a letter, and let
\m{\mathsf{w}\colon\fS\to\WW}, \m{\mathsf{t}_0\colon\fS\to\fS} be maps with
$s-\mathsf{w}(s)=\cE\,\mathsf{t}_0(s)$ for all \m{s\in\fS}. Let $v_1,\dots,v_m\in\constE{l}{A}{i}(L)$,
\m{c\in\constE{l}{A}{i}(L)}, \m{n\ge0} an integer, \m{s_1,\dots,s_m\in\fS}, $f$ a polynomial with integer
coefficients, with $f_n$ its coefficient of degree~$n$, and $o_1\colon\mathbb{Z}\to\mathbb{Z}$ a function with
\m{o_1(u)\le A} for all~$u$. Suppose that
$\sum_as_av_a-f(\cE)\,c\in\wExp{L}{i}{\shStd{l}{o_1}{\cdot}{\cdot}}$. Then the element
\[d=\sum_a\mathsf{w}\bigl(\mathsf{t}_0^{\,n}(s_a)\bigr)v_a-f_n\,c\]
lies in \m{\constE{l}{A}{i}(L)} and in $\wExp{L}{i}{\shStd{l}{o_3}{\cdot}{\cdot}}$, where \m{o_3(u)=0} if
\m{o_1(u)\le n} and \m{o_3(u)=A} otherwise.
\end{lemma}
\leanline{\lean{coeffmemconstantsModE}\sep the formal statement introduces~$d$ by a local definition inside the
statement}
\begin{proof}
The module \m{\constE{l}{A}{i}(L)} is stable under~$\WW$, whose elements are fixed by~$\Gamma'$; so $d$ lies in
it. By induction on~$n$. Write $s_a=\mathsf{w}(s_a)+\cE\mathsf{t}_0(s_a)$ and \m{f=f_0+Xf_1}:
\[\sum_as_av_a-f(\cE)c=d_0+\cE\,y,\qquad d_0=\sum_a\mathsf{w}(s_a)v_a-f_0c,\quad
y=\sum_a\mathsf{t}_0(s_a)v_a-f_1(\cE)c .\]
Here $y$ has the exponents~$\std_l$, so \m{\cE y} has them raised by~$1$ at \m{\nu=0}, and $d_0$ has the
exponents $\shStd{l}{\min(1,o_1)}{\cdot}{\cdot}$. Its Galois differences have the exponent~$A$ at \m{\nu=0}.
By Input~\ref{q:deep}, $d_0$ has at \m{\nu=0} the exponent~$A$ where \m{o_1(u)\ge1} and \m{\min(A,0)=0}
elsewhere: this is the case \m{n=0}. For the step, \m{d_0} has the exponents
\m{\shStd{l}{o_1}{\cdot}{\cdot}}, hence so has \m{\cE y}, and $y$ has the exponents
\m{\shStd{l}{o_1-1}{\cdot}{\cdot}}. The induction hypothesis for \m{(\mathsf{t}_0(s_a))_a}, $f_1$,
\m{o_1-1} and~$n$ gives the claim for \m{n+1}.
\end{proof}

\begin{lemma}[projections that are constant modulo a power of~$\cE$]\label{l:projconst}
Let $L$ be a lattice datum, \m{l>0} and \m{A\ge0} integers, $i$ a letter and \m{x\in\WC\otimes L}. Suppose that
$x$ is a point over~$F_0$ of \m{\varphi^{l}\Minf(L)} and that \m{g\mu(x)-\mu(x)\in\Fil^{A}\BdR} for all
\m{\mu\in\Per(L)} and \m{g\in\Gamma'}. Then \m{\pr_i(x)\in\constE{l}{A}{i}(L)}.
\end{lemma}
\leanline{\lean{letterProjmemconstantsModE}}
\begin{proof}
\m{\pr_i(x)\in\locLP{l}{i}(L)} by Input~\ref{q:B}. Its Galois differences have the exponents~$\std_l$; at
\m{\nu=0}, where \m{\std_l=0} because \m{l>0}, the value of a period~$\mu$ on \m{g\pr_i(x)-\pr_i(x)} is
$(g-1)(\mu\circ\pr_i)(x)\in\Fil^{A}\BdR$. Input~\ref{q:exp} concludes.
\end{proof}

\begin{proof}[Proof of Theorem D]
\emph{A multiple in the Frobenius pullback.} By Lemma~\ref{l:first} there is
$y\in\ptsOver{S}{c}{i}{0}{\raisedE{-\rho}{\cdot}{\cdot}}{\raisedE{-\rho}{\cdot}{\cdot}}$. As the chart is
unramified, $\cE_a=z^{p^{a}}+p=(\cE-p)^{p^{a}}+p$. Let
\[f=X^{\rho}\prod_{0\le a<l}\bigl((X-p)^{p^{a}}+p\bigr)^{\gamma},\]
a monic polynomial of degree~$N$ with integer coefficients; then
$f(\cE)y=\cE^{\rho}\prod_{a<l}\cE_a^{\gamma}\,y$ has the exponents \m{\gamma-u} at \m{\nu<l} and \m{-u} at
\m{\nu\ge l} (Input~\ref{q:exp}). Since the jumps at these~$\nu$ are at most~$\gamma$, it has the
exponents~$\std_l$, so it is a point over~$F_0$ of \m{\varphi^{l}\Minf(\extL{S})} (Input~\ref{q:B}).

\emph{Spanning.} By Input~\ref{j:frame}, $p^{4\delta}f(\cE)y=\sum_ar_av_a$ with \m{r_a\in\fS} and
points~$v_a$ whose values are $\Gamma'$-invariant modulo \m{\Fil^{A}\BdR}. Projecting to the letter~$i$ and
using Lemma~\ref{l:projconst}: $p^{4\delta}f(\cE)y=\sum_ar_a\pr_i(v_a)$ with
$\pr_i(v_a)\in\constE{l}{A}{i}(\extL{S})$.

\emph{The coefficient of degree~$N$.} Choose \m{\mathsf{w},\mathsf{t}_0} as in Lemma~\ref{l:coeff}
(Input~\ref{j:mod}) and put $d=\sum_a\mathsf{w}(\mathsf{t}_0^{\,N}(r_a))\pr_i(v_a)$. Apply the lemma with
\m{c=0}, \m{f=0} and \m{o_1(u)=\min(A,\gamma-u)}: $d$ lies in \m{\constE{l}{A}{i}(\extL{S})} and has at
\m{\nu=0} the exponent~$A$ where \m{\gamma-u>N}, and~$0$ elsewhere. Apply it also in \m{\trivL{S}}, to the
images of the \m{\pr_i(v_a)}, to the element \m{\constAt{S}{i}{c}}, the polynomial \m{p^{4\delta}f} and the
constant function~$A$: the hypothesis holds because $y$ lies over \m{\constAt{S}{i}{c}}, and as $f$ is monic
of degree \m{N<A} the conclusion is that
\[a_0=(1\otimes\projC{S})(d)-p^{4\delta}\constAt{S}{i}{c}\]
has the exponents \m{\shStd{l}{o_2}{\cdot}{\cdot}} in \m{\trivL{S}}, \m{o_2} the constant~$A$. So
$\cE^{-A}a_0\in\locLP{l}{i}(\trivL{S})$.

\emph{The point.} From $\Kills{S}{i^{p^{-l}}}{\kappa}{0}$ we get $p^{\kappa}\in\liftId{S}{0}{l}{i}$, so by
Input~\ref{q:kill} there is \m{v\in\locLP{l}{i}(\extL{S})} over \m{p^{\kappa}\cE^{-A}a_0}. The element
\m{x=p^{\kappa}d-\cE^{A}v} lies in \m{\Mloc(\extL{S})}, over $p^{\kappa+4\delta}\constAt{S}{i}{c}$. At
\m{\nu=0} its exponent is at least \m{\min(A,-u-N)}: the term \m{\cE^{A}v} has~$A$, and $p^{\kappa}d$ has~$A$
or, where \m{\gamma-u\le N}, the exponent \m{0\ge-u-N}. Elsewhere the exponents are~$\std_l$. The Galois
differences of both terms have the exponents \m{\shStd{l}{o_2}{\cdot}{\cdot}}.
\end{proof}

\sect{Where the results are used}{s:used}
Part~3 of the notes on this unit applies Theorem~B with cuts chosen on the Frobenius orbits of the letters and
windows around them, Theorem~C for lattices whose weights lie in a low and a far window, Theorem~D and
Proposition~\ref{p:auer} for classes with few low weights, and the lemmas on points for the cokernel killer of
tame types. Its two results, a length bound and an exponent, are the inputs of the bound for Selmer groups
at~$\ell$ in the unit on Selmer groups.

\sect{Notes on the formal text}{s:notes}
\emph{Arguments as the text gives them.} The arguments were read from the Lean text and the census of the
workspace, and were not printed by Lean. The census reads a bound variable as an explicit argument in the
definition of the cost of a map (a letter) and in the statement of Proposition~\ref{p:reach}; the note states
what the Lean text binds. The property of Definition~\ref{d:deep} has eleven arguments, which the note writes
out in words. Two statements contain a local definition before their conclusion, which a program that cuts a
statement at its first definition sign misses: Lemma~\ref{l:third} and Lemma~\ref{l:coeff}; both were read in the
source in full. The hypotheses of Proposition~\ref{p:safe} and Lemma~\ref{l:safediff} are section hypotheses of
the formal text, included in these two declarations by name.

\emph{Differences from the natural statements.} A scaled cup product is asked to bound for every topology
on~$\Gamma$ that makes $j_\Gamma$ an open embedding; no topology is part of the data. In Theorem~B the function
\m{\mathsf{l}} is an argument together with the hypothesis that it gives the nearest cut; so every Frobenius
orbit of letters must contain a cut. The families $z'$ of depths of the third points occur in the hypotheses of
Theorem~B and are not used in its proof. In Definition~\ref{d:deep} the integer~$r$ and the property that
\m{i^{p^{-r}}} is a cut locate the third point used with clause~(b). In Theorem~C the same pair
\m{(\kappa,\rho)} kills the three cokernels at every letter. In Proposition~\ref{p:auer} the formal text takes
\m{hn+1} elements indexed from~$0$.

\emph{Arithmetic.} In the proofs of Theorems B and~C the formal text settles the comparisons of exponents by a
decision procedure for linear arithmetic. The note gives them by hand: in Theorem~B as the reduction of each
safety condition to a clause of Definition~\ref{d:deep}, in Theorem~C as the explicit inequalities of the two
alternatives. In the second alternative of Theorem~C the hypothesis \m{p\ge5} enters only through the cokernel
killer of Part~1, and \m{\eta\ge1} only at \m{\nu=0} for the letters with \m{\varsigma(i)=1}.

\emph{The background.} The period ring is Colmez's~$B_{\max}$. For the Artin--Schreier--Witt input the second
reader of the workspace finds the Frobenius surjectivity in the literature and no source for the full statement
used here; for Auerbach's lemma a determinant-maximization extension; for Tate's theorem a weaker form. No place
in the literature was confirmed for the frame theorem and for the stability of the class of lattice data.

\subhead{How this note was written.} From the reader's dossier of the unit: the statements of all members and the
structured proofs of the declarations treated here. The two Lean files were read in full. The proofs of
Theorem~A and of Theorem~B, which the reader had marked as least certain, were read line by line in the Lean
text: the identities for \m{D_g} in the proof of Theorem~A were recomputed by hand, and the four non-trivial
safety conditions of the walk, with the lemma and the third point used in each, were matched with the four
tactic blocks of the formal proof. The inequalities in the proof of Theorem~C, in both alternatives, were worked out by
hand for this note; the formal proof leaves them to automation, so the account of where the windows, the
excluded interval and the hypothesis \m{\eta\ge1} enter is ours. The proof of Theorem~D and of Proposition~\ref{p:auer} follow the formal proofs. The program
check of the unit's dossier covered \UnitSteps{} leaf steps of \UnitProofs{} proofs.

\head{Appendix A. Concordance with the Lean artifact}
{\fontsize{8.6}{11.6}\selectfont\renewcommand{\leanfont}{\ttfamily\fontsize{7.7}{10}\selectfont\color{leancol}}
\setlength{\LTleft}{0pt}\setlength{\LTright}{\fill}\setlength{\tabcolsep}{0pt}
\begin{longtable}{@{}p{30mm}p{56mm}p{43mm}r@{}}
\textbf{In this note} & \textbf{Lean declaration} & \textbf{File} & \textbf{Lines}\\\midrule\endhead
\ConcordanceRows
\end{longtable}}

Two statements as the artifact has them:
\quoted{ofisSafe}
\quoted{HasDeepPoint}

\emph{Legend.} As in Part~1; in addition \texttt{P} or \texttt{h} is the pairing~$\Pi$ and \texttt{P.m} is~$M$;
\texttt{smallPart} is~$\mathfrak{s}_M$; \texttt{zmodPowToUnits} is~$\mathsf{e}_M$; \texttt{thirdPairing}
is~$\mathsf{p}_\Pi$; \texttt{toZp} is~$\operatorname{tr}$; \texttt{cut}, \texttt{Q} and \texttt{l} of the walk
are $\mathsf{cut}$, $\mathsf{Q}$ and~$\mathsf{l}$; the integers \texttt{s}, \texttt{a} that are exponents
of~$\cE$ are~$\rho$, \texttt{b} is~$\beta$, \texttt{ε} is~$\varsigma$, and \texttt{g} of the long gap
is~$\gamma$.

\head{Notation}
\begingroup\fontsize{8.6}{10.9}\selectfont\raggedright
\setlength{\columnsep}{9mm}\setlength{\parskip}{0pt}
\begin{multicols}{2}
\nt{$\ell$}{the prime of Chapter~VIII of the handout; it is~$p$ here}
\nt{$p$, $\Qp$, $\Zp$, $\bar{\mathbb{Q}}_p$, $C$, $\mathcal{O}_C$, $\GQp$}{the prime, the $p$-adic numbers and integers, an algebraic closure, its completion, the integers, the Galois group}
\nt{$\Ainf$, $\WC$, $W(\mathfrak{m}^\flat)$, $\mathcal{O}_C^\flat$, $C^\flat$, $\mathfrak{m}^\flat$}{Witt vectors of the tilt, of its fraction field, of its maximal ideal}
\nt{$\varphi$, $\theta$, $\varepsilon$, $\tau$, $[\,\cdot\,]$}{Frobenius, Fontaine's map, a system of roots of unity, \m{[\varepsilon^{1/p}]-1}, the Teichmüller lift}
\nt{$\BdR$, $B^{+}_{\mathrm{dR}}$, $\Fil^{n}\BdR$, $B_{\max}$}{period rings and the filtration}
\nt{$\Gamma$, $\Gamma'$, $j_\Gamma$, $\chi$, $\Delta$}{the group, the subgroup of the chart, the map to~$\GQp$, the cyclotomic character, an open subgroup}
\nt{$\fc$, $\fS$, $\WW$, $h$, $X$, $z$, $\cE$, $\cE_\nu$, $e_{\fc}$, $\fd$}{the chart, its ring, \m{W(\mathbb{F}_{p^h})}, the residue degree, the variable, the coordinate, the Eisenstein element and its translates, the index, the different exponent}
\nt{$F$, $F_\nu$, $\operatorname{co}_0$}{the field fixed by~$\Gamma'$, the value fields, the retraction \m{\Ainf\to\fS}}
\nt{$\cO$, $\omega$, $\cO(1)$}{coefficients, the root of unity, the Tate twist}
\nt{$i$, $j$, $j'$, $i^-$, $j^-$, $j^{\circ}$}{letters; \m{i^-=i^{p^{-1}}}; the nearest cut before~$j$, the cut from which a point is moved}
\nt{$L$, $L'$, $A$, $C$, $T$, $T'$, $Y$, $Y'$}{lattice data and lattices; $A$ is also an integer}
\nt{$U_L(i)$, $m_L(i,u)$}{the jump set and the multiplicity of a jump at a letter}
\nt{$\pr_i$, $\Per(L)$, $\Per_{i,u}(L)$}{projector and periods}
\nt{$\Minf(L)$, $\Mlocinf(L)$, $\Mloc(L)$, $\ret_L$}{Kisin modules and the retraction}
\nt{\m{\wExp{L}{i}{o}}, \m{\locLP{l}{i}(L)}, \m{\locKP{Q}}, \m{\invMod{Q}}}{exponent modules; elements of \m{\Mloc}; invariants modulo~$Q$}
\nt{$\std_l$, \m{\cost{\nu}{n}}, $\impE{o^{\mathrm{lo}}}{o^{\mathrm{hi}}}{\nu}{u}$, \m{\frobX{l}{n}{o}{\nu}{u}}}{standard exponents, cost, improved and shifted exponents}
\nt{\m{\wPair{b}{y}{x}}}{the extension of a bilinear map~$b$}
\nt{$S$, $S'$, $Y_S$, $C_S$, $\alpha_S$}{sides and their parts}
\nt{\m{\extL{S}}, \m{\trivL{S}}, \m{\projC{S}}, \m{\projY{S}}, \m{\inclY{S}}, \m{\projCloc{S}}}{extension, trivial datum, projections, inclusion}
\nt{\m{\constAt{S}{i}{c}}, \m{\liftId{S}{r}{l}{i}}, \m{\Kills{S}{i}{\kappa}{\rho}}}{constants, lifting ideal, cokernel killer}
\nt{\m{\raisedE{n}{\nu}{u}}, \m{\twoWin{n}{c}}, \m{\NoMult{p}{n}{\beta}}, $\JumpWin{\eta}{\varsigma}{\beta}{W}{U}$}{raised exponents and windows}
\nt{$\Pi$, $M$, $\psi$, $\operatorname{tr}$, $\bar{\operatorname{tr}}$, $c_\Pi$, $\langle\ ,\ \rangle$, $\mathsf{p}_\Pi$}{a pairing and its parts; the third pairing}
\nt{$\mathbb{M}_M$, $\mathfrak{s}_M$, $\mathsf{e}_M$}{the module of values, its small part, the embedding of \m{\mathbb{Z}/p^{M}}}
\nt{\m{\tensL{S}{S'}}, \m{\tateL{\cO}}, \m{\third{\Pi}}, \m{\pairW{\Pi}}}{tensor datum, Tate datum, third side, extended pairing}
\nt{\m{\Safe{\Pi}{e}{w}{y}}, \m{\CupB{\Pi}{e}{c}{c'}}}{a safe pair; the scaled cup product bounds}
\nt{\m{\kerX{S}{i}{o}}, \m{\deepE{S}{i}{o}}, \m{\Reaches{S}{e}{\mathsf{A}}}, \m{\ptsOver{S}{c}{i}{e}{o_1}{o'}}}{kernel with exponents, deep elements, cost of a map, points}
\nt{\m{\raiseX{o}{\nu}{u}{\nu'}{v}}, \m{\shStd{l}{o_1}{\nu}{u}}, \m{\constE{l}{A}{i}(L)}}{an exponent raised at one place; shifted standard exponents; constants modulo~$\cE^{A}$}
\nt{$\mathsf{cut}$, $\mathsf{Q}$, $\mathsf{q}$, $\mathsf{l}(i)$}{properties of letters, a proposition, the distance to the nearest cut}
\nt{$o$, $o'$, $o''$, $o_1,\dots,o_5$, $o_+$, $o^{\mathrm{lo}}$, $o^{\mathrm{hi}}$, $d$, $d'$, $d_j$, $d'_j$, $D$, $z(i)$, $z'(i)$}{exponent functions and families of them; $o_1$, $o_2$, $o_3$ are also functions of~$u$ alone; $z$ is also the coordinate}
\nt{$\nu$, $\nu'$, $\nu_0$, $l$, $u$, $u'$, $u_3$, $v$}{indices, levels, jumps; $v$ is also an element}
\nt{$e$, $e'$, $e_1$, $e_2$, $e_3$, $E$, $\kappa$, $\rho$, $\epsilon$, $H$, $N$, $W$, $\eta$, $\beta$, $\gamma$, $\delta$, $n$, $n_0$, $m$, $q$, $r$, $r'$, $r_j$, $r'_j$, $a$, $b$, $k$}{integers; $n$ is also a pair or a point, $q$ also an element, $r$ also a power of~$p$, $k$ also a cochain or an index, $b$ also a bilinear map}
\nt{$\varsigma$}{a sign}
\nt{$g$, $h$, $t_k$}{elements of~$\Gamma$; $h$ is also the residue degree}
\nt{$x$, $x'$, $x_1$, $x_2$, $x_j$, $x'_j$, $x_k$, $w$, $w_0$, $y$, $y'$, $y_0$, $y_c$, $y_k$, $y_{kb}$, $n_j$, $n_1$, $v_j$, $u_j$, $v_a$, $a_0$, $c$, $c'$, $c_0$, $c_k$, $s$, $s_a$, $r_a$, $f$, $f_n$, $f_0$, $f_1$, $t$, $t'$, $\lambda$, $\pi$}{elements of modules, lattices and rings; $f$ is also an equivariant map}
\nt{$X^{\top}_j$, $X^{\bot}_j$, $X^{\mathsf{Q}}_j$, $X'^{\top}_j$, $X'^{\bot}_j$ and their sums}{the points of the walk}
\nt{$\xi$, $\xi'$, $\xi_w$, $\mathsf{a}_g$, $\mathsf{a}'_g$, $\mathsf{B}$, $D_g$, $y_0$, $k(g)$, $\mathsf{c}$}{the objects of the proof of Theorem~A; $\mathsf{B}$ is also a map in Proposition~\ref{p:reach}}
\nt{$\mu$, $\mu'$, $\mu_C$, $\mu_m$, $\Lambda$, $a_c$, $V$}{periods, a linear map, elements of~$F$ and of~$\BdR$}
\nt{$\mathsf{A}$, $\mathsf{A}_0$, $I$, $I_0$}{additive maps and sets of steps}
\nt{$\mathsf{m}$, $\mathsf{t}$}{scalar multiplication and the tensor map}
\nt{$t_{km}$, $\bar t_{km}$, $a_{kb}$, $\bar a_{kb}$, $w_{kb}$, $s_b$, $\mathsf{U}$, $\mathsf{z}_b$, $\mathsf{v}$}{the matrices of the proof of Proposition~\ref{p:auer} and the objects of Auerbach's lemma}
\nt{$\mathsf{w}$, $\mathsf{t}_0$}{remainder and quotient of the division by~$\cE$}
\nt{$M$, $N$, $Q$}{submodules; $M$ is also the exponent of the pairing, $N$ also an integer}
\end{multicols}
\endgroup

\end{document}
